
\documentclass[12pt]{article}

\usepackage{amsmath, amsthm, amssymb, mathtools}
\usepackage{geometry}
\usepackage{hyperref}
\usepackage{enumitem}
\usepackage{microtype}
\usepackage{xcolor}

\geometry{margin=1.25in}

\newtheorem{theorem}{Theorem}[section]
\newtheorem{proposition}[theorem]{Proposition}
\newtheorem{lemma}[theorem]{Lemma}
\newtheorem{corollary}[theorem]{Corollary}
\newtheorem{conjecture}[theorem]{Conjecture}
\theoremstyle{definition}
\newtheorem{definition}[theorem]{Definition}
\newtheorem{example}[theorem]{Example}
\theoremstyle{remark}
\newtheorem{remark}[theorem]{Remark}
\newtheorem{warning}[theorem]{Warning}

\DeclareMathOperator{\supp}{supp}
\DeclareMathOperator{\type}{type}
\DeclareMathOperator{\Stab}{Stab}
\newcommand{\R}{\mathbb{R}}
\newcommand{\C}{\mathbb{C}}
\newcommand{\Z}{\mathbb{Z}}
\newcommand{\N}{\mathbb{N}}
\newcommand{\T}{\mathcal{T}}
\newcommand{\PW}{\mathrm{PW}}
\newcommand{\Sph}{\mathbb{S}}
\newcommand{\wh}{\widehat}
\newcommand{\abs}[1]{\left|#1\right|}
\newcommand{\norm}[1]{\left\|#1\right\|}
\newcommand{\ip}[2]{\left\langle #1,\, #2\right\rangle}
\newcommand{\re}{\operatorname{Re}}

% Status markers
\newcommand{\proven}{\textbf{\color{teal}[Proved]}}
\newcommand{\conditional}{\textbf{\color{orange}[Conditional]}}
\newcommand{\open}{\textbf{\color{red}[Open]}}
\newcommand{\heuristic}{\textbf{\color{purple}[Heuristic]}}

\title{\textbf{Nonlinear Convolution Equations, Spectral Rigidity,\\
and the Sharp Restriction Constant on $S^1$}\\[0.5em]
\large A Compactness Route to Finiteness of Critical Values}
\author{}
\date{\today}

\begin{document}

\maketitle
\tableofcontents
\bigskip

% =======================================================================
\section{Overview and Strategy}
% =======================================================================

The ultimate target of this programme is the sharp Fourier extension
constant on the unit circle $S^1\subset\R^2$:
\[
  \mathcal R
  :=
  \sup_{0\neq f\in L^2(S^1)}
  \frac{\norm{\wh{f\sigma}}_{L^6(\R^2)}}{\norm{f}_{L^2(S^1)}},
\]
where $\sigma$ is arclength measure on $S^1$.  Every extremiser satisfies a
nonlinear convolution equation on the sphere that fits into a general
framework studied in the first part of this note.

\subsection{The main strategy}

After normalising $\norm{f}_{L^2}=1$ and writing $u=\wh{f\sigma}$, an
extremiser satisfies the Euler--Lagrange equation
\[
  u = c\bigl((|u|^4u)\ast\wh\sigma\bigr),
  \qquad c=\Lambda^{-1},\qquad \Lambda=\norm{u}_{L^6}^6.
\]
The strategy is:
\begin{enumerate}[leftmargin=2em]
  \item \textbf{Spectral support forcing.}  Show that $\wh u$ is supported
    on~$S^1$, reducing from all of~$\R^2$ to a density
    $\rho\in L^2(S^1)$.  \proven
  \item \textbf{Fredholm structure.}  Prove that the linearisation of the
    Euler--Lagrange map at any critical point is Fredholm of index zero,
    so that solutions are locally isolated modulo symmetry.  \proven
  \item \textbf{Compactness modulo symmetry.}  Show that the set of normalised
    critical points is sequentially compact after quotienting by the symmetry
    group $G=U(1)\times SO(2)\ltimes\R^2$.  \open
  \item \textbf{Nondegeneracy.}  At each critical point, verify that the kernel
    of the linearisation equals the tangent space to the symmetry orbit.
    At the constant solution, this reduces to finitely many Bessel integral
    inequalities.  \conditional
  \item \textbf{Finiteness and identification.}  Conclude that the critical
    values $\Lambda_j$ are finitely many, and identify the sharp constant as
    $\mathcal R = \max_j \Lambda_j^{1/6}$.  \conditional
\end{enumerate}

\subsection{Status summary}
\renewcommand{\arraystretch}{1.4}
\begin{center}
\begin{tabular}{lll}
\hline
\textbf{Component} & \textbf{Status} & \textbf{Key input} \\
\hline
Spectral support forcing (any $n$) &
  \proven & Paley--Wiener + support calculus \\
Compactness of $T$ on $\PW$ (general) &
  \proven & Rellich for PW spaces \\
Compactness of $K_{u_0}$ on $L^2(S^1)$ &
  \proven & Truncation + Hilbert--Schmidt \\
Fredholm structure ($I-K$ index 0) &
  \proven & Riesz--Schauder \\
1D case ($S^0$) & \proven & Explicit algebraic solution \\
Nondegeneracy at constant solution &
  \conditional & Bessel integrals $\mathcal J_n$ \\
Existence of maximiser &
  \proven & Profile decomposition + strict convexity \\
Compactness modulo $G$ (all critical pts) &
  \open & Ruling out multi-profile splitting \\
Finiteness of critical values &
  \conditional & Items 3 + 4 above \\
Higher-dimensional generalisations &
  \open & Incidence-fibre smoothing \\
\hline
\end{tabular}
\end{center}

% =======================================================================
% PART I: GENERAL THEORY
% =======================================================================

\part{General Theory}

% =======================================================================
\section{Setup and Notation}
\label{sec:setup}
% =======================================================================

We work throughout with \emph{entire functions of finite exponential type}.
An entire function $u:\C\to\C$ has \emph{exponential type} $\tau\geq 0$ if
$\abs{u(z)}\leq Ce^{\tau\abs{z}}$ for some $C>0$ and all $z\in\C$.
The \emph{Paley--Wiener--Schwartz theorem} identifies this class:
$u$ is entire of type $\leq\tau$ with $u|_{\R}\in L^2(\R)$ if and only if
$\wh u\in L^2$ is supported in $[-\tau,\tau]$.
We write $\PW_\tau$ for the corresponding Hilbert space.

Fourier conventions:
\[
  \wh u(\xi)=\int_\R u(x)e^{-ix\xi}\,dx,
  \qquad
  u(x)=\frac{1}{2\pi}\int_\R \wh u(\xi)e^{ix\xi}\,d\xi,
\]
so $\widehat{u\ast v}=\wh u\cdot\wh v$ and
$\widehat{u\cdot v}=\tfrac{1}{2\pi}\wh u\ast\wh v$.

Throughout, $k\geq 1$ is a fixed positive integer, $c\in\C\setminus\{0\}$
is a fixed constant, and $f$ is an entire function of exponential type
$\tau_f>0$ whose restriction to~$\R$ lies in $L^2(\R)$.

% =======================================================================
\section{The Two Convolution Equations}
\label{sec:two-eqns}
% =======================================================================

\subsection{The equation $u^k=c\,(u\ast f)$: spectral expansion kills solutions}

In Fourier space the equation reads
$\wh u^{\ast k}(\xi)=c\,\wh u(\xi)\wh f(\xi)$.

\begin{proposition}[Type obstruction]  \label{prop:type-obstruction}
  Let $u\in\PW_\tau$ with $\tau>0$ satisfy $u^k=c(u\ast f)$ for $k\geq 2$.
  Then $u\equiv 0$.
\end{proposition}

\begin{proof}
  The Titchmarsh convolution theorem gives
  $\supp(\wh u^{\ast k})\subseteq[-k\tau,k\tau]$, while the right side is
  supported in $[-\tau,\tau]$.  For $k\geq 2$ this forces $\tau=0$, so $u$
  is a polynomial; degree comparison then forces $u$ to be constant, and the
  only constants satisfying the equation are the $(k-1)$-th roots of
  $c\wh f(0)$ together with $0$.
\end{proof}

\begin{corollary}
  For $k\geq 2$ the solution set of $u^k=c(u\ast f)$ in the Paley--Wiener
  class is a finite discrete set of at most $k$ constants.
\end{corollary}

\subsection{The equation $u=c\,(u^k\ast f)$: compact fixed-point structure}
\label{sec:reversed-eqn}

We now turn to the equation
\begin{equation}  \label{eq:main}
  u = c\,(u^k\ast f).
\end{equation}
In Fourier space:
\begin{equation}  \label{eq:main-fourier}
  \wh u(\xi) = c\,\wh u^{\ast k}(\xi)\cdot\wh f(\xi).
\end{equation}

\begin{proposition}[Type self-consistency]
  Any solution $u\in\PW_\tau$ of~\eqref{eq:main} satisfies $\tau\leq\tau_f$.
\end{proposition}

\begin{proof}
  The right side of~\eqref{eq:main-fourier} is supported in
  $\supp(\wh u^{\ast k})\cap\supp(\wh f)\subseteq[-\tau_f,\tau_f]$.
\end{proof}

This is the key contrast: the equation $u=c(u^k\ast f)$ does \emph{not}
force $\tau=0$; solutions of positive exponential type exist as long as
$\tau\leq\tau_f$.

\begin{proposition}[Compactness of $T$ on Paley--Wiener spaces]
  \label{prop:T-compact-PW}
  Define $T:\PW_{\tau_f}\to\PW_{\tau_f}$ by $T(v)=c(v^k\ast f)$.  Then
  $T$ maps bounded sets in $\PW_\tau$ (for any $\tau>\tau_f$) to precompact
  sets in $\PW_\tau$.  In particular, $T:\PW_\tau\to\PW_\tau$ is a compact
  operator.
\end{proposition}

\begin{proof}
  Since $v\in\PW_{\tau_f}$ implies $v^k\in\PW_{k\tau_f}$, the function
  $v^k\ast f$ has Fourier support in
  $\supp(\wh{v^k})\cap\supp(\wh f)\subseteq[-\tau_f,\tau_f]$.
  Hence $T$ maps $\PW_\tau$ into $\PW_{\tau_f}$, and the inclusion
  $\PW_{\tau_f}\hookrightarrow\PW_\tau$ is compact by the Rellich-type
  theorem for Paley--Wiener spaces.
\end{proof}

\subsection{Fredholm structure and local manifold theorem}

At a solution $u_0$ of~\eqref{eq:main}, the Fr\'echet derivative of
$G(u):=u-c(u^k\ast f)$ is
\[
  DG(u_0)[h] = h - ck(u_0^{k-1}h\ast f) =: (I-K)[h],
\]
where $K:h\mapsto ck(u_0^{k-1}h\ast f)$ is compact by the same argument.
By Riesz--Schauder, $I-K$ is \textbf{Fredholm of index zero}.

\begin{theorem}[Local manifold structure]
  \label{thm:local-mfd}
  Let $u_0$ be a solution of~\eqref{eq:main} at which $DG(u_0)=I-K$ is
  invertible.  Then $u_0$ is an isolated solution.  More generally, the
  solution set near $u_0$ is a smooth manifold of dimension
  $\dim\ker(I-K)$, which is finite.
\end{theorem}

\begin{remark}[Direct contraction argument]
  For $u,v$ in a ball of radius $R$ in $\PW_{\tau_f}$, Young's inequality
  gives
  $\norm{T(u)-T(v)} \leq \abs{c}kR^{k-1}\norm{f}_{L^1}\norm{u-v}$.
  This is a contraction whenever $R<(\abs{c}k\norm{f}_{L^1})^{-1/(k-1)}$,
  giving local uniqueness without spectral theory.
\end{remark}

\subsection{Comparison of the two equations}

\renewcommand{\arraystretch}{1.5}
\begin{table}[h]
\centering
\begin{tabular}{lll}
\hline
& $u^k=c(u\ast f)$ & $u=c(u^k\ast f)$ \\
\hline
Fourier identity & $\wh u^{\ast k}=c\wh u\wh f$ &
  $\wh u=c\wh u^{\ast k}\cdot\wh f$ \\
Type constraint ($k\geq 2$) & Forces $\tau=0$ (Titchmarsh) &
  Allows $\tau\leq\tau_f$ \\
Operator structure & No compactness & $T$ compact \\
Linearisation & Not Fredholm in general &
  $I-K$, Fredholm index $0$ \\
Solutions ($k\geq 2$, generic) &
  Finitely many constants & Isolated; Leray--Schauder applies \\
\hline
\end{tabular}
\caption{Analytic contrast between the two convolution equations.}
\end{table}

% =======================================================================
\section{Spectral Support Forcing}
\label{sec:spectral-forcing}
% =======================================================================

We specialise to $f=\wh\sigma$, where $\sigma$ is the surface measure on
$S^{n-1}\subset\R^n$.  By the Paley--Wiener theorem, $f$ is entire of
exponential type~$1$, and $\wh f=(2\pi)^n\sigma$ as a distribution.

\begin{proposition}[Spectral support forcing]
  \label{prop:support-forcing}
  If $u$ is entire of exponential type satisfying $u=c(u^k\ast\wh\sigma)$,
  then $\wh u$ is a distribution supported on~$S^{n-1}$.
  The same conclusion holds for the equation $u=c(|u|^{2(k-1)}u\ast\wh\sigma)$.
\end{proposition}

\begin{proof}
  From~\eqref{eq:main-fourier}, $\wh u=c\,\wh u^{\ast k}\cdot(2\pi)^n\sigma$.
  The right side is supported on $\supp\sigma=S^{n-1}$.

  For the $|u|^{2(k-1)}u$ variant: writing $|u|^{2(k-1)}u=u^k\bar u^{k-1}$
  in Fourier space involves $(k-1)$-fold convolutions of $\wh{\bar u}$, where
  $\wh{\bar u}(\xi)=\overline{\wh u(-\xi)}$.  Since $S^{n-1}$ is symmetric
  under $\xi\mapsto-\xi$, the support of all convolved terms lies in the
  sumset of copies of $\pm S^{n-1}$, and multiplication by $\sigma$ restricts
  back to $S^{n-1}$.
\end{proof}

\begin{remark}
  In dimension $n\geq 2$, the conclusion $\supp\wh u\subseteq S^{n-1}$ means
  $u$ is a \emph{Herglotz wave function}, satisfying $(-\Delta-1)u=0$.
\end{remark}

\subsection{Reduction to a density equation on the sphere}

Write $\wh u=\rho\,d\sigma$ for $\rho\in L^2(S^{n-1})$, so
\[
  u(x) = \int_{S^{n-1}}\rho(\omega)e^{ix\cdot\omega}\,d\sigma(\omega).
\]
The equation $u=c(u^k\ast\wh\sigma)$ then reduces to the
\emph{spherical fixed-point equation}
\begin{equation}  \label{eq:sphere-eq}
  \rho(\omega) = c\int_{V_k(\omega)}\rho(\omega_1)\cdots\rho(\omega_k)
    \,d\sigma_{\mathrm{fiber}}(\omega_1,\ldots,\omega_k),
  \qquad \sigma\text{-a.e.}\ \omega\in S^{n-1},
\end{equation}
where $V_k(\omega)=\{(\omega_1,\ldots,\omega_k)\in(S^{n-1})^k:
\omega_1+\cdots+\omega_k=\omega\}$ is the incidence fibre.

\begin{warning}[Incidence fibre $\neq$ spherical convolution]
  \label{warn:not-convolution}
  The operator $T(\rho)(\omega)=\int_{V_k(\omega)}\rho^{\otimes k}\,
  d\sigma_\mathrm{fiber}$ is \emph{not} the same as convolution by the
  spherical measure.  In particular, the Funk--Hecke theorem does not
  diagonalise~$T$ in spherical harmonics.  Any reduction to finitely many
  harmonic coefficients requires separate justification.
\end{warning}

% =======================================================================
\section{The One-Dimensional Case: Complete Solution}
\label{sec:1D}
% =======================================================================

For $n=1$, $S^0=\{-1,+1\}$, $\sigma=\delta_{-1}+\delta_{+1}$, and
$f(x)=\wh\sigma(x)=2\cos x$.  Proposition~\ref{prop:support-forcing} forces
$\wh u=a\delta_{-1}+b\delta_{+1}$, so $u(x)=ae^{-ix}+be^{ix}$.

\subsection{Algebraic reduction}

Expanding $u^k$ and computing $u^k\ast f$ picks out frequencies
$\pm 1$, which exist only when $k$ is \textbf{odd}.

\begin{proposition}[Parity obstruction]
  If $k$ is even, the equation $u=c(u^k\ast f)$ with $f=2\cos$ has only
  $u\equiv 0$.
\end{proposition}

For $k$ odd, setting $\beta_k:=\binom{k}{(k-1)/2}$ and comparing coefficients
of $e^{\mp ix}$ gives the single constraint
\begin{equation}  \label{eq:constraint}
  ab = \lambda_k,
  \qquad
  \lambda_k := \left(\frac{1}{2\pi c\beta_k}\right)^{2/(k-1)}.
\end{equation}

\begin{theorem}[Complete solution in 1D]  \label{thm:main-1d}
  \proven\quad
  Let $k\geq 3$ be odd, $c\in\C\setminus\{0\}$, $f=2\cos$.
  The solution set of $u=c(u^k\ast f)$ in the Paley--Wiener class is
  \[
    \{0\}
    \;\cup\;
    \bigsqcup_{j=0}^{(k-3)/2}
      \{(a,b)\in\C^2:ab=\omega^j\lambda_k\},
    \qquad \omega=e^{4\pi i/(k-1)}.
  \]
  This is a finite union of $(k-1)/2$ smooth one-dimensional complex
  manifolds (each $\cong\C^*$), plus the origin.
  Each component is exactly the translation orbit of any single non-trivial
  solution on it: the translation $u(x)\mapsto u(x-t)$ acts by
  $(a,b)\mapsto(e^{-it}a,e^{it}b)$, preserving~$ab$.
\end{theorem}

\begin{example}[$k=3$]
  $\beta_3=3$.  Constraint: $ab=(6\pi c)^{-1}$.  One component, the
  translation orbit of $u(x)=(6\pi c)^{-1/2}\cos x$.
\end{example}

\begin{example}[$k=5$]
  $\beta_5=10$.  Constraint: $ab=(20\pi c)^{-1/2}$.  Two components
  (the scaling group $\mu_2=\{1,-1\}$ generates a second hyperbola).
\end{example}

\begin{remark}[Real solutions]
  For real-valued solutions $b=\bar a$, so $ab=|a|^2\geq 0$.  Only the
  hyperbola with $\lambda_k\in\R_{>0}$ contributes; the real solution
  manifold is the circle $|a|=|\lambda_k|^{1/2}$, parametrised by
  translations.
\end{remark}

% =======================================================================
% PART II: THE S^1 RESTRICTION PROBLEM
% =======================================================================

\part{The $S^1$ Restriction Problem}

% =======================================================================
\section{The Sharp Restriction Constant and Its Euler--Lagrange Equation}
\label{sec:restriction}
% =======================================================================

\subsection{Setup}

The sharp Fourier extension constant on $S^1\subset\R^2$ is
\[
  \mathcal R
  :=
  \sup_{\substack{f\in L^2(S^1)\\ f\neq 0}}
  \frac{\norm{Ef}_{L^6(\R^2)}}{\norm{f}_{L^2(S^1)}},
  \qquad
  (Ef)(x):=\int_{S^1}f(\omega)e^{ix\cdot\omega}\,d\sigma(\omega),
\]
where $\sigma$ is the arclength measure on $S^1$ (total mass $2\pi$).
Every extension $u=Ef$ satisfies $(-\Delta-1)u=0$.

\subsection{The Euler--Lagrange equation}

Normalise $\norm{f}_{L^2(S^1)}=1$.  A critical point of
$f\mapsto\norm{Ef}_{L^6}^6$ on the unit sphere satisfies
\begin{equation}  \label{eq:EL}
  E^*(|u|^4u) = \Lambda f,
  \qquad u=Ef,\qquad \Lambda=\norm{u}_{L^6}^6,
\end{equation}
where $E^*:L^{6/5}(\R^2)\to L^2(S^1)$ is the adjoint restriction:
$(E^*g)(\omega)=\int_{\R^2}g(x)e^{-ix\cdot\omega}\,dx$ restricted to
$\omega\in S^1$.  Equivalently,
\begin{equation}  \label{eq:EL-u}
  u = c\bigl((|u|^4u)\ast\wh\sigma\bigr),
  \qquad c=\Lambda^{-1}.
\end{equation}

\begin{warning}[The nonlinearity is $|u|^4u$, not $u^5$]
  The restriction problem gives $|u|^4u=u^3\bar u^2$, which is
  degree~$5$ but \emph{not} holomorphic.  The linearisation is
  \emph{real-linear}, not complex-linear: the Jacobian couples the Fourier
  mode~$n$ to mode~$-n$.  This distinction is essential for the spectral
  analysis in Section~\ref{sec:nondeg}.
\end{warning}

\subsection{The symmetry group}

The natural symmetry group is
\[
  G = U(1)\times SO(2)\ltimes\R^2,
  \qquad \dim G=4.
\]
It acts on densities by
\[
  (e^{i\alpha},R,x_0)\cdot f(\omega)
  = e^{i\alpha}e^{-ix_0\cdot\omega}f(R^{-1}\omega).
\]
Phase, rotation, and modulation/translation all preserve $\norm{f}_{L^2}$
and $\norm{Ef}_{L^6}$, hence preserve~$\Lambda$.  Every critical point
generates a $G$-orbit of critical points with the same~$\Lambda$.

\begin{definition}[Shape]
  A \emph{shape} is an equivalence class of normalised critical points modulo
  the $G$-action.  Two critical points with the same shape have the same
  critical value~$\Lambda$.
\end{definition}

\subsection{What ``finding all critical values'' means}

The search problem is: determine the finite set
\[
  \mathcal Z
  :=
  \{(f,\Lambda):E^*(|Ef|^4Ef)=\Lambda f,\;\norm{f}_2=1\}/G.
\]
If $\mathcal Z$ is finite, the sharp restriction constant is
\[
  \mathcal R = \max_{[f,\Lambda]\in\mathcal Z}\Lambda^{1/6}
  = c_{\min}^{-1/6},
\]
where $c_{\min}$ is the smallest positive $c_j=\Lambda_j^{-1}$.
The normalisation $\norm{f}_2=1$ is essential: without it, $\alpha u$ solves
the equation with $c$ replaced by $c\alpha^{-4}$, preventing discreteness.

\subsection{Positive representatives for extremisers}

The exponent $6$ gives a useful reduction for the \emph{extremiser}
search.

\begin{lemma}[Even-exponent positivity reduction]
  \label{lem:positive-reduction}
  For every $f\in L^2(S^1)$,
  \[
    \norm{Ef}_{L^6(\R^2)} \leq \norm{E|f|}_{L^6(\R^2)}.
  \]
  Consequently, if $f$ is an extremiser, then $|f|$ is also an extremiser.
  In particular, to identify the sharp constant it is enough to search
  among nonlinear eigenfunctions with a nonnegative representative.
\end{lemma}

\begin{proof}
  By Plancherel and the fact that $6=2\cdot 3$, up to the fixed Fourier
  normalisation constant,
  \[
    \norm{Ef}_{L^6}^6
    = \norm{(f\sigma)*(f\sigma)*(f\sigma)}_{L^2(\R^2)}^2 .
  \]
  The convolution density satisfies the pointwise estimate
  \[
    \abs{(f\sigma)*(f\sigma)*(f\sigma)}
    \leq (|f|\sigma)*(|f|\sigma)*(|f|\sigma)
  \]
  almost everywhere.  Squaring and integrating, then applying Plancherel
  in the reverse direction, gives the asserted inequality.

  Since $\norm{|f|}_{L^2(S^1)}=\norm{f}_{L^2(S^1)}$, replacing an
  extremiser by its modulus preserves the quotient and hence produces
  another extremiser.
\end{proof}

\begin{remark}
  Lemma~\ref{lem:positive-reduction} is an extremiser reduction, not a
  classification of all critical points.  A non-maximal critical point need
  not have a nonnegative representative.  Thus the full compactness theorem
  for every critical sequence remains stronger than what is needed merely
  to compute~$\mathcal R$.
\end{remark}

% =======================================================================
\section{Fredholm Structure of the Linearisation}
\label{sec:fredholm}
% =======================================================================

\subsection{The extension operator $E$ is not compact}

One might hope that $E:L^2(S^1)\to L^6(\R^2)$ is compact, which would
immediately give compactness of the nonlinear map $T(f)=E^*(|Ef|^4Ef)$.
This is \textbf{false}.

\begin{proposition}[$E$ is not compact]
  \label{prop:E-not-compact}
  \proven\quad
  The extension operator $E:L^2(S^1)\to L^6(\R^2)$ is bounded
  (Stein--Tomas) but not compact.
\end{proposition}

\begin{proof}
  Fix $f\in L^2(S^1)$ with $\norm{f}_2=1$ and $\norm{Ef}_{L^6}>0$.
  For $n\in\N$, define $f_n(\omega):=e^{-ine_1\cdot\omega}f(\omega)$
  where $e_1=(1,0)$.  Then $\norm{f_n}_2=1$ and
  $Ef_n(x)=Ef(x-ne_1)=u(x-ne_1)$, a translate of $u=Ef$.

  For $n\neq m$ with $|n-m|$ large, the supports of $|u(\cdot-ne_1)|$ and
  $|u(\cdot-me_1)|$ in any $L^6$-significant region are essentially disjoint
  (since $u(x)\to 0$ as $|x|\to\infty$ by Riemann--Lebesgue), so
  \[
    \norm{Ef_n-Ef_m}_{L^6}
    = \norm{u(\cdot-ne_1)-u(\cdot-me_1)}_{L^6}
    \;\xrightarrow{|n-m|\to\infty}\;
    2^{1/6}\norm{u}_{L^6}>0.
  \]
  Hence $\{Ef_n\}$ has no Cauchy subsequence.
\end{proof}

The obstruction is the noncompact translation symmetry: high-frequency
modulations on $S^1$ shift the extension arbitrarily far in $x$-space.
This means compactness of the critical set requires quotienting by~$G$;
it cannot be deduced from a single operator bound.

\subsection{The linearised operator $K_{u_0}$ is compact}

Despite the non-compactness of $E$, the \emph{specific} operator appearing
in the linearisation is compact, thanks to the decay of the solution $u_0$.

\begin{definition}
  At a critical point $f_0$ with $u_0=Ef_0$, define the linearised operator
  $K=K_{u_0}:L^2(S^1)\to L^2(S^1)$ by
  \[
    Kh := E^*\bigl(DN(u_0)[Eh]\bigr),
  \]
  where $DN(u_0)[v]=3|u_0|^4v+2|u_0|^2u_0^2\bar v$ is the (real-linear)
  Fr\'echet derivative of $N(u)=|u|^4u$.
\end{definition}

\begin{theorem}[Compactness of $K_{u_0}$]
  \label{thm:K-compact}
  \proven\quad
  For any $u_0=Ef_0$ with $f_0\in L^2(S^1)$, the operator
  $K_{u_0}:L^2(S^1)\to L^2(S^1)$ is compact.
\end{theorem}

\begin{proof}
  We prove compactness for the ``main term''
  $K_1 h:=E^*(|u_0|^4 Eh)$; the argument for
  $K_2 h:=E^*(|u_0|^2 u_0^2\overline{Eh})$ is identical,
  and $K=3K_1+2K_2$.

  \medskip\noindent
  \textbf{Step 1: Truncation.}
  For $R>0$ let $\chi_R$ denote the indicator of $B(0,R)\subset\R^2$ and
  set $K_1^R h:=E^*(\chi_R\,|u_0|^4\,Eh)$.  Then
  \[
    \norm{K_1-K_1^R}_{L^2\to L^2}
    \leq
    \norm{E^*}_{L^{6/5}\to L^2}
    \;\norm{(1-\chi_R)|u_0|^4}_{L^{3/2}}
    \;\norm{E}_{L^2\to L^6}.
  \]
  Here $\norm{E}$ and $\norm{E^*}$ are the (finite) Stein--Tomas constants,
  and
  \[
    \norm{(1-\chi_R)|u_0|^4}_{L^{3/2}}^{3/2}
    = \int_{|x|>R}|u_0|^6\,dx
    \;\xrightarrow{R\to\infty}\; 0,
  \]
  since $u_0\in L^6(\R^2)$.  Hence $\norm{K_1-K_1^R}_{\mathrm{op}}\to 0$.

  \medskip\noindent
  \textbf{Step 2: $K_1^R$ is Hilbert--Schmidt.}
  The integral kernel of $K_1^R$ on $S^1\times S^1$ is
  \[
    \mathcal K_R(\omega,\omega')
    = \int_{B(0,R)}|u_0(x)|^4\,e^{ix\cdot(\omega'-\omega)}\,dx.
  \]
  Since $|u_0|^4\chi_R\in L^1(\R^2)\cap L^\infty(\R^2)$ (bounded and
  compactly supported), its Fourier transform is continuous on $\R^2$.
  Therefore $\mathcal K_R$ is continuous on the compact set $S^1\times S^1$,
  hence bounded, and in particular belongs to $L^2(S^1\times S^1)$.
  Thus $K_1^R$ is Hilbert--Schmidt, hence compact.

  \medskip\noindent
  \textbf{Step 3: Conclusion.}
  $K_1=\lim_{R\to\infty}K_1^R$ in operator norm, where each $K_1^R$ is
  compact.  Hence $K_1$ is compact.
\end{proof}

\begin{corollary}[Fredholm structure]  \label{cor:fredholm}
  \proven\quad
  At any critical point $f_0$ of the restriction functional with
  $\norm{f_0}_2=1$ and $\Lambda_0=\norm{Ef_0}_6^6>0$, the operator
  \[
    L := K_{u_0} - \Lambda_0 I
  \]
  has the form $L=K-\Lambda_0 I$ with $K$ compact.  Therefore
  $L:L^2(S^1)\to L^2(S^1)$ is Fredholm of index zero.
  In particular, $\ker L$ is finite-dimensional, and $L$ is invertible
  if and only if $\ker L=\{0\}$.
\end{corollary}

\begin{remark}[Self-adjointness]
  With respect to the real inner product
  $\re\ip{f}{g}_{L^2(S^1)}$, the operator $K$ is self-adjoint.
  Indeed, $DN(u_0)$ is self-adjoint on $L^2(\R^2,\re\ip{\cdot}{\cdot})$
  (it is the Hessian of the real functional $u\mapsto\tfrac{1}{6}|u|^6$),
  and $E^*$ is the adjoint of $E$.  Consequently, the eigenvalues of $L$
  are real.
\end{remark}

% =======================================================================
\section{Nondegeneracy at the Constant Solution}
\label{sec:nondeg}
% =======================================================================

\subsection{Setup}

The constant function $f_0=(2\pi)^{-1/2}$ has $\norm{f_0}_{L^2(S^1)}=1$.
Its extension is the radial function
\[
  u_0(x) = Ef_0(x) = (2\pi)^{1/2}J_0(|x|),
\]
where $J_0$ is the Bessel function of order zero.  Set $A:=(2\pi)^{1/2}$, so
$u_0(r)=AJ_0(r)$.  The critical value is
\[
  \Lambda = \norm{u_0}_{L^6}^6
  = A^6\cdot 2\pi\int_0^\infty J_0(r)^6\,r\,dr
  = 16\pi^4\int_0^\infty J_0(r)^6\,r\,dr.
\]

The stabiliser of $f_0$ under $G$ is $\Stab_G(f_0)=SO(2)$ (rotations of
$S^1$ fix the constant function).  The orbit has dimension
$\dim G-\dim\Stab=4-1=3$, generated by phase ($U(1)$) and translations
($\R^2$).

\subsection{Fourier decomposition of the linearised operator}

Since $u_0$ is radial, the linearised operator
$K:L^2(S^1)\to L^2(S^1)$ preserves each pair of Fourier modes
$\{e^{in\theta},e^{-in\theta}\}$.  We use the Fourier expansion
\[
  h(\theta)=\sum_{m\in\Z}c_m e^{im\theta},
  \qquad
  Eh(r,\phi)=\sum_{m\in\Z} 2\pi\,i^{|m|}c_m\,J_{|m|}(r)\,e^{im\phi},
\]
where $(r,\phi)$ are polar coordinates in $\R^2$.

\begin{proposition}[Eigenvalue structure]
  \label{prop:eigenvalues}
  \proven\quad
  Define the Bessel integrals
  \begin{equation}  \label{eq:Jn-def}
    \mathcal J_n := \int_0^\infty J_0(r)^4\,J_n(r)^2\,r\,dr,
    \qquad n\geq 0,
  \end{equation}
  and $\mathcal I_n:=16\pi^4\mathcal J_n$, so that $\Lambda=\mathcal I_0$.
  Then, on the sector $\{c_ne^{in\theta}+c_{-n}e^{-in\theta}\}$ with
  $n\geq 1$, the operator $K=E^*DN(u_0)E$ acts as a $4\times 4$
  real-linear map with eigenvalues
  \[
    5\mathcal I_n\quad(\text{multiplicity }2)
    \qquad\text{and}\qquad
    \mathcal I_n\quad(\text{multiplicity }2).
  \]
  On the sector $n=0$, the eigenvalues are $5\mathcal I_0$ (for real
  perturbations) and $\mathcal I_0$ (for imaginary perturbations), each
  with multiplicity~$1$.
\end{proposition}

\begin{proof}
  Write $v=Eh$ where $h=c_ne^{in\theta}+c_{-n}e^{-in\theta}$.  Then
  $v=2\pi i^n J_n(r)[c_ne^{in\phi}+c_{-n}e^{-in\phi}]$
  (using $E(e^{im\theta})=2\pi i^{|m|}J_{|m|}(r)e^{im\phi}$).

  Since $u_0=AJ_0(r)$ is real and radial,
  $DN(u_0)[v]=3u_0^4 v+2u_0^4\bar v$. The $e^{in\phi}$ component of
  $DN(u_0)[v]$ is $2\pi A^4J_0^4J_n[3i^nc_n+2(-i)^n\bar c_{-n}]$,
  and the $e^{-in\phi}$ component is
  $2\pi A^4J_0^4J_n[3i^nc_{-n}+2(-i)^n\bar c_n]$.

  Applying $E^*$ and using the identity
  $\langle E^*g,e^{im\theta}\rangle=2\pi(-i)^{|m|}\int_0^\infty
  g_m(r)J_{|m|}(r)\,r\,dr$
  (where $g_m$ is the $m$-th angular Fourier coefficient of~$g$),
  together with $(-i)^ni^n=1$ and $(-i)^{2n}=(-1)^n$, one computes:
  \begin{align}
    (Kh)_n &= \mathcal I_n\bigl[3c_n+2(-1)^n\bar c_{-n}\bigr],
    \label{eq:Kn}\\
    (Kh)_{-n} &= \mathcal I_n\bigl[3c_{-n}+2(-1)^n\bar c_n\bigr].
    \label{eq:K-n}
  \end{align}
  Writing $c_n=p+iq$, $c_{-n}=s+it$ (real and imaginary parts), the system
  decouples into a \emph{real part} and an \emph{imaginary part}, each a
  $2\times 2$ real system.  For both parities of~$n$, both $2\times 2$
  matrices have the same eigenvalues: $5$ and $1$ (with the eigenvectors
  exchanging between the real and imaginary blocks as $(-1)^n$ changes sign).
  Hence the eigenvalues of~$K$ on the sector are $5\mathcal I_n$ and
  $\mathcal I_n$, each with multiplicity~$2$.
  The $n=0$ case is analogous with multiplicities halved.
\end{proof}

\subsection{Identification of the symmetry kernel}

\begin{proposition}[Symmetry directions]
  \label{prop:symmetry-kernel}
  \proven\quad
  The kernel of $L=K-\Lambda I$ on $L^2(S^1)$ includes the following
  directions, with the indicated eigenvalue identities.
  \begin{enumerate}[leftmargin=2em]
    \item \textbf{Phase} ($n=0$, imaginary):
      $h=if_0$.  This lies in the $\mathcal I_0$-eigenspace of~$K$.
      Since $\Lambda=\mathcal I_0$, the $L$-eigenvalue is $0$.
    \item \textbf{Translations} ($n=1$):
      $h=-i(\hat x_0\cdot\omega)f_0$ for $\hat x_0\in\R^2$.
      This lies in the $5\mathcal I_1$-eigenspace of~$K$.
      Since translations are an exact symmetry, $5\mathcal I_1=\Lambda$.
    \item \textbf{No rotation direction}:
      $\partial_\theta f_0=0$ since $f_0$ is constant.
  \end{enumerate}
  In particular, $\mathcal I_1=\Lambda/5$, i.e.,
  \begin{equation}  \label{eq:bessel-identity}
    5\int_0^\infty J_0(r)^4\,J_1(r)^2\,r\,dr
    = \int_0^\infty J_0(r)^6\,r\,dr.
  \end{equation}
  This is a Bessel function identity forced by translation invariance.
\end{proposition}

\begin{proof}
  (i) is verified by direct computation: $E(if_0)=iu_0$,
  $DN(u_0)[iu_0]=3u_0^4(iu_0)+2u_0^4(-iu_0)=iu_0^5$, and
  $E^*(iu_0^5)=iE^*(u_0^5)=i\Lambda f_0=\Lambda(if_0)$.

  (ii) The translation $f\mapsto e^{-ix_0\cdot\omega}f$ gives
  $h=-i(\hat x_0\cdot\omega)f_0$ at first order.  The corresponding
  extension perturbation is $Eh=-\hat x_0\cdot\nabla u_0$
  (verified by differentiating $E(e^{-ix_0\cdot\omega}f_0)(x)=u_0(x-x_0)$
  in $x_0$).  Differentiating the Euler--Lagrange equation
  $E^*(|u_0(\cdot-x_0)|^4 u_0(\cdot-x_0))=\Lambda e^{-ix_0\cdot\omega}f_0$
  in $x_0$ at $x_0=0$ gives $Kh=\Lambda h$, confirming that $h$ is in
  $\ker L$.

  From the eigenvalue structure, $h=-i(\hat x_0\cdot\omega)f_0$ has
  Fourier support on $n=\pm 1$ with $c_{-1}=-\bar c_1$.  Substituting
  into~\eqref{eq:Kn}--\eqref{eq:K-n} shows this lies in the
  $5\mathcal I_1$-eigenspace, whence $5\mathcal I_1=\Lambda$.

  (iii) $f_0$ is constant, so $\partial_\theta f_0=0$. $\square$
\end{proof}

\begin{remark}[Euler identity]
  The identity $Kf_0=5\Lambda f_0$ (in the $n=0$ real sector) follows from
  the degree-$5$ homogeneity of $N$: $DN(u)[u]=5N(u)$ for all $u$.
  Hence $Kf_0=E^*(DN(u_0)[u_0])=5E^*(N(u_0))=5\Lambda f_0$.
\end{remark}

\subsection{The nondegeneracy criterion}

The eigenvalues of $L=K-\Lambda I$ on the $n$-th sector are
$5\mathcal I_n-\Lambda$ and $\mathcal I_n-\Lambda$.  At $n=0$: $0$ (phase)
and $4\Lambda$ (amplitude).  At $n=1$: $0$ (translations, multiplicity~2) and
$-4\Lambda/5$ (multiplicity~2).

\begin{theorem}[Nondegeneracy reduces to Bessel integrals]
  \label{thm:nondeg-criterion}
  \conditional\quad
  The constant solution $f_0=(2\pi)^{-1/2}$ is nondegenerate modulo
  $G$ (i.e., $\ker L=T_{f_0}(G\cdot f_0)$ on the tangent space to
  $\norm{f}_2=1$) if and only if
  \begin{equation}  \label{eq:nondeg-cond}
    \boxed{
    \mathcal I_n\neq\Lambda
    \quad\text{and}\quad
    5\mathcal I_n\neq\Lambda
    \qquad\text{for all } n\geq 2.}
  \end{equation}
  Equivalently, in terms of the Bessel integrals~\eqref{eq:Jn-def}:
  \[
    \mathcal J_n\neq\mathcal J_0
    \quad\text{and}\quad
    5\mathcal J_n\neq\mathcal J_0
    \qquad\text{for all } n\geq 2.
  \]
\end{theorem}

\begin{proof}
  By Proposition~\ref{prop:eigenvalues}, the eigenvalues of $L$ on the
  $n$-th sector ($n\geq 2$) are $5\mathcal I_n-\Lambda$ and
  $\mathcal I_n-\Lambda$, each with multiplicity~$2$.  These are both
  nonzero if and only if~\eqref{eq:nondeg-cond} holds.  The symmetry
  directions exhaust the kernel at $n=0$ and $n=1$ by
  Proposition~\ref{prop:symmetry-kernel}.
\end{proof}

\subsection{Asymptotic analysis of $\mathcal J_n$}

\begin{lemma}[Decay of $\mathcal J_n$]
  \label{lem:Jn-decay}
  \proven\quad
  $\mathcal J_n\to 0$ as $n\to\infty$.  More precisely,
  $\mathcal J_n=O(n^{-1})$.
\end{lemma}

\begin{proof}
  For $r<n(1-\varepsilon)$, the Bessel function $J_n(r)$ is exponentially
  small.  For $r\geq n$, the uniform bound $|J_0(r)|\leq 1$ and the
  asymptotic $|J_\nu(r)|\leq Cr^{-1/2}$ (valid for $r\geq\nu$ and all
  $\nu\geq 0$) give
  \[
    \int_n^\infty J_0(r)^4\,J_n(r)^2\,r\,dr
    \leq C\int_n^\infty r^{-2}\cdot r^{-1}\cdot r\,dr
    = C\int_n^\infty r^{-2}\,dr = C/n.  \qedhere
  \]
\end{proof}

Since $\mathcal J_n\to 0$ while $\mathcal J_0>0$ and
$\mathcal J_0/5>0$, the conditions~\eqref{eq:nondeg-cond} hold for all
sufficiently large~$n$.  Only finitely many values of~$n$ need to be
checked.

\begin{lemma}[Parity structure of the asymptotics]
  \label{lem:parity}
  For moderate~$n$ (before the $J_n$-support retreats past the main
  lobe of $J_0$), the large-$r$ average of $J_0(r)^4J_n(r)^2$ depends
  on the parity of~$n$.  Using the asymptotics
  $J_\nu(r)\sim\sqrt{2/\pi r}\cos(r-\nu\pi/2-\pi/4)$ for $r\gg\nu$:
  \begin{itemize}
    \item \textbf{Even~$n$:}
      $\cos(r-n\pi/2-\pi/4)=\pm\cos(r-\pi/4)$, so
      $\langle J_0^4J_n^2\rangle_{\mathrm{avg}}
      \sim(2/\pi r)^3\cdot\tfrac{5}{16}$,
      the same coefficient as~$\langle J_0^6\rangle$.
    \item \textbf{Odd~$n$:}
      $\cos(r-n\pi/2-\pi/4)=\pm\sin(r-\pi/4)$, so
      $\langle J_0^4J_n^2\rangle_{\mathrm{avg}}
      \sim(2/\pi r)^3\cdot\tfrac{1}{16}$,
      which is $\tfrac{1}{5}$ of the even case.
  \end{itemize}
  Consequently, for moderate even~$n$, $\mathcal J_n$ is close to
  $\mathcal J_0$ (deficiency only from the small-$r$ region where
  $J_n(r)\ll J_0(r)$), while for moderate odd~$n$, $\mathcal J_n$ is
  close to $\mathcal J_0/5=\mathcal J_1$.
\end{lemma}

\begin{remark}[Which conditions are tight?]
  The identity $\mathcal J_1=\mathcal J_0/5$ (from translation invariance)
  shows that the condition $5\mathcal J_n=\mathcal J_0$ \emph{is} achieved
  at $n=1$.  For even~$n$, the condition $\mathcal J_n=\mathcal J_0$ is
  never achieved ($\mathcal J_n<\mathcal J_0$ due to the small-$r$
  deficiency), but the margin may be small for $n=2$.  The most delicate
  check is therefore whether $5\mathcal J_n=\mathcal J_0$ for some
  \emph{even} $n\geq 2$ (where $\mathcal J_n$ is close to $\mathcal J_0$,
  hence $5\mathcal J_n$ is close to $5\mathcal J_0\gg\Lambda$, making
  this condition safe), and whether $\mathcal J_n=\mathcal J_0$ for some
  \emph{odd} $n\geq 3$ (where $\mathcal J_n\approx\mathcal J_0/5$, well
  below $\mathcal J_0$, also safe).

  \textbf{Upshot:} the asymptotic structure makes the nondegeneracy
  conditions~\eqref{eq:nondeg-cond} very plausible, but does not close the
  argument.  A rigorous verification requires evaluating $\mathcal J_n$ for
  $n=2,3,\ldots,n_*$ (where $n_*$ is the threshold beyond which the
  $O(1/n)$ bound suffices) and confirming the two inequalities at each.
  This is a finite numerical computation that can be made rigorous with
  interval arithmetic.
\end{remark}

\subsection{Nondegeneracy at non-constant solutions}
\label{sec:nondeg-nonconst}

If there exist non-constant critical points of the restriction functional,
their nondegeneracy must be established separately.  At a non-constant
$f_0$, the extension $u_0=Ef_0$ is not radial, so $K_{u_0}$ does not
decompose into Fourier sectors.  Two approaches are available:

\paragraph{Direct spectral analysis.}
If $u_0$ has a discrete symmetry (e.g., if $f_0$ is a real-valued
even function on $S^1$), the problem decomposes into finitely many
symmetry sectors, each finite-dimensional.

\paragraph{Generic transversality (Sard--Smale).}
Introduce a family of $G$-equivariant perturbations $F_p(\rho)=\rho-cT_p(\rho)$
and show that the parameter derivatives span the cokernel of $DF_p(\rho)$
at every zero.  This gives nondegeneracy for a generic parameter~$p$.
To cover the original (unperturbed) problem, one must either verify that
the spherical measure is a regular parameter, or continue the generic
solution set to the original parameter without bifurcation.

\begin{remark}[Varying $c$ alone is insufficient]
  As noted in Warning~\ref{warn:not-convolution}, varying only the scalar
  $c$ provides a one-parameter family, which is generically too small for
  Sard--Smale (the cokernel at a degenerate zero may have dimension
  $>1$).  A richer family of perturbations is needed.
\end{remark}

% =======================================================================
\section{Compactness Modulo Symmetry}
\label{sec:compactness}
% =======================================================================

\subsection{Regularity bootstrap}

\begin{proposition}[Smoothness and gauge-fixed uniform bounds]
  \label{prop:bootstrap}
  \proven\quad
  If $f\in L^2(S^1)$ satisfies $E^*(|Ef|^4Ef)=\Lambda f$ with
  $\Lambda>0$, then $f\in C^\infty(S^1)$.  Moreover, for each
  $k\geq 0$ there exists $C_k=C_k(\varepsilon)$ such that
  $\norm{f}_{C^k(S^1)}\leq C_k$ for all normalised critical points
  with $\Lambda\geq\varepsilon$ in a fixed translation gauge; in
  particular, this applies in the nonnegative slice used in
  Lemma~\ref{lem:positive-gauge}.
\end{proposition}

\begin{proof}[Proof sketch]
  Write $u=Ef$, $g=|u|^4u$.  Since $\wh u=f\,d\sigma$ is supported on
  $S^1$, the Fourier transform $\wh g$ is supported in the Minkowski sum
  $S^1+S^1+S^1+S^1+S^1=\overline{B(0,5)}$ (using $|u|^4u=u^3\bar u^2$
  and $\wh{\bar u}$ supported on $-S^1=S^1$).  Thus $g$ is
  \emph{frequency-localised} to $B(0,5)$.

  The Euler--Lagrange equation gives $f=\Lambda^{-1}\wh g\big|_{S^1}$,
  so $f$ inherits the regularity of~$\wh g$ restricted to the compact
  curve $S^1\subset B(0,5)$.  The $n$-th Fourier coefficient of $f$
  on $S^1$ is
  \begin{equation}\label{eq:fourier-bootstrap}
    \wh f(n)
    = \Lambda^{-1}\cdot 2\pi(-i)^{|n|}
      \int_0^\infty g_n(r)\,J_{|n|}(r)\,r\,dr,
  \end{equation}
  where $g_n(r)$ is the $n$-th angular Fourier coefficient of~$g$.
  For $|n|\gg 1$, $J_{|n|}(r)$ is exponentially small for $r<|n|/2$
  (Debye bounds), and the contribution from $r>|n|/2$ is controlled by
  the $L^{6/5}$ tail of~$g$ via H\"older and the Bessel asymptotics
  $J_\nu(r)=O(r^{-1/2})$ for $r>\nu$.  An inductive argument (improved
  decay of $\wh f(n)$ feeds back into improved regularity of $f$ and
  hence of $u$ and~$g$) yields $|\wh f(n)|=O(|n|^{-N})$ for all~$N$.
  The constants are obtained from the $L^2$ normalisation, the lower
  bound $\Lambda\geq\varepsilon$, and the Euler--Lagrange equation, by
  this energy/bootstrap scheme.

  The translation-gauge qualification is necessary.  If
  $f_x(\omega)=e^{-ix\cdot\omega}f(\omega)$, then $f_x$ is again a
  normalised critical point with the same~$\Lambda$, while
  $\norm{f_x}_{C^1}$ grows like~$|x|$ in general.  Once the translation
  parameter is fixed, or once one works in the nonnegative slice below,
  this obstruction is absent.
\end{proof}

\begin{corollary}[Uniform pointwise decay in a fixed gauge]
  \label{cor:uniform-decay}
  For all normalised critical points in a fixed translation gauge with
  $\Lambda\geq\varepsilon$,
  \[
    |Ef(x)|\leq C(\varepsilon)\,(1+|x|)^{-1/2}
    \qquad\text{for all }x\in\R^2.
  \]
  Equivalently, for an arbitrary translated representative the same
  estimate holds around its translation centre.
\end{corollary}
\begin{proof}
  By the bootstrap, $\norm{f}_{C^1}\leq C(\varepsilon)$.  The standard
  stationary-phase estimate for extension from a $C^\infty$ curve in
  $\R^2$ with $C^1$ density gives $|Ef(x)|=O(|x|^{-1/2})$
  for $|x|\geq 1$, and
  $|Ef(x)|\leq(2\pi)^{1/2}\norm{f}_{L^2}\leq(2\pi)^{1/2}$ everywhere.
\end{proof}

\subsection{Profile decomposition for the circle extension}

The extension $E:L^2(S^1)\to L^6(\R^2)$ is not compact
(Proposition~\ref{prop:E-not-compact}), and the sole obstruction is
the translation symmetry.  The following decomposition, analogous to
results of Merle--Vega and Shao for the paraboloid, captures this
defect of compactness completely.

\begin{theorem}[Profile decomposition]
  \label{thm:profile-decomp}
  \conditional\quad
  Let $f_j\in L^2(S^1)$ with $\norm{f_j}_2\leq C$.  After passing
  to a subsequence, there exist profiles $\phi^k\in L^2(S^1)$
  ($k=1,2,\ldots$) and group elements
  $g_j^k=(e^{i\alpha_j^k},R_j^k,x_j^k)\in G$ satisfying the
  \emph{divergence condition}
  $|x_j^k-x_j^l|\to\infty$ for all $k\neq l$,
  such that for each $K\geq 1$,
  \begin{equation}\label{eq:profile-decomp}
    f_j = \sum_{k=1}^K(g_j^k)^{-1}\cdot\phi^k + r_j^K,
  \end{equation}
  with the decoupling identities
  \begin{align}
    \norm{f_j}_2^2
    &= \textstyle\sum_{k=1}^K\norm{\phi^k}_2^2
      + \norm{r_j^K}_2^2 + o(1),
    \label{eq:L2-decouple}\\
    \norm{Ef_j}_6^6
    &= \textstyle\sum_{k=1}^K\norm{E\phi^k}_6^6
      + \norm{Er_j^K}_6^6 + o(1),
    \label{eq:L6-decouple}
  \end{align}
  and $\norm{Er_j^K}_{L^6}\to 0$ as $K\to\infty$ (after further
  subsequence extraction at each stage).
\end{theorem}

The $L^2$ decoupling~\eqref{eq:L2-decouple} is elementary:
for $k\neq l$,
$\ip{e^{-ix_j^k\cdot\omega}\phi^k}{e^{-ix_j^l\cdot\omega}\phi^l}
=\int_{S^1}\phi^k\overline{\phi^l}\,
e^{i(x_j^l-x_j^k)\cdot\omega}\,d\sigma\to 0$
by the Riemann--Lebesgue lemma on $S^1$.

The $L^6$ decoupling~\eqref{eq:L6-decouple} is deeper.  Since
$E(e^{-ix_j^k\cdot\omega}\phi^k)(x)=(E\phi^k)(x-x_j^k)$, distinct
profiles produce extensions centred at diverging locations.  A
Br\'ezis--Lieb argument using the Herglotz wave decay
$|E\phi^k(x)|=O(|x|^{-1/2})$ (valid once $\phi^k$ is smooth;
for $L^2$ profiles, one approximates and uses density) then gives the
$L^6$ asymptotic additivity.

The \emph{construction} of the decomposition proceeds iteratively.
The key input is an inverse Strichartz inequality for $S^1$:

\begin{lemma}[Inverse Strichartz for $S^1$]
  \label{lem:inv-strichartz}
  If $\norm{f_j}_2=1$ and $\norm{Ef_j}_{L^6}\geq\delta>0$, then
  there exist $x_j\in\R^2$ such that the re-centred sequence
  $\tilde f_j(\omega):=e^{ix_j\cdot\omega}f_j(\omega)$ has a
  subsequence converging weakly to some $\phi\neq 0$ in $L^2(S^1)$.
\end{lemma}

\begin{proof}
  Since $\norm{Ef_j}_6^6\geq\delta^6$, a covering argument gives
  $x_j\in\R^2$ with $\int_{B(x_j,1)}|Ef_j|^6\geq c\delta^6$.
  Set $\tilde u_j(x):=Ef_j(x+x_j)=E\tilde f_j(x)$.  Then
  $\int_{B(0,1)}|\tilde u_j|^6\geq c\delta^6$.

  Every Herglotz wave satisfies $(\Delta+1)u=0$ and
  $\norm{u}_{L^\infty}\leq(2\pi)^{1/2}\norm{f}_2=(2\pi)^{1/2}$.
  By interior Schauder estimates for $\Delta u+u=0$, the family
  $\{\tilde u_j\}$ is equicontinuous on $\overline{B(0,1)}$.
  Arzel\`a--Ascoli extracts a uniformly convergent subsequence
  $\tilde u_j\to u_\infty$ on $\overline{B(0,1)}$.

  Meanwhile, $\tilde f_j\rightharpoonup\phi$ (Banach--Alaoglu,
  after subsequence).  Weak convergence gives
  $E\tilde f_j(x)\to E\phi(x)$ pointwise for every~$x$
  (since $\omega\mapsto e^{ix\cdot\omega}\in L^2(S^1)$).  So
  $u_\infty=E\phi$.  From
  $\int_{B(0,1)}|u_\infty|^6\geq c\delta^6>0$, we get
  $\phi\neq 0$.
\end{proof}

\begin{remark}[Simplification from compact frequency support]
  For the paraboloid $\{(\xi,|\xi|^2)\}$, profile decompositions must
  account for parabolic rescalings (a second family of symmetries).
  For $S^1$, the frequency set is compact with no scaling symmetry,
  so translations are the sole defect of compactness.  Moreover, the
  inverse Strichartz lemma follows from elliptic regularity for
  $\Delta u+u=0$, rather than from bilinear restriction estimates.
\end{remark}

\subsection{Profiles of critical sequences inherit the equation}

The following result is the key structural input: if the original
sequence consists of critical points, then each profile in the
decomposition is itself a critical point.

\begin{proposition}[Each profile is a critical point]
  \label{prop:profile-critical}
  \proven\quad
  Let $(f_j,\Lambda_j)$ be normalised critical points with
  $\Lambda_j\to\Lambda_\infty>0$, and let $\phi^1,\phi^2,\ldots$
  be the profiles from Theorem~\ref{thm:profile-decomp} (re-indexed
  so that $x_j^1=0$).  Then each profile satisfies
  \begin{equation}\label{eq:profile-EL}
    E^*(|E\phi^k|^4E\phi^k) = \Lambda_\infty\phi^k.
  \end{equation}
  In particular, $\psi^k:=\phi^k/\norm{\phi^k}_2$ is a normalised
  critical point with eigenvalue
  $\widetilde\Lambda_k=\Lambda_\infty/\norm{\phi^k}_2^4$.
\end{proposition}

\begin{proof}
  We prove~\eqref{eq:profile-EL} for $k=1$; the case $k\geq 2$
  follows by applying the translation $y_j^k:=x_j^k-x_j^1$
  (which diverges) and repeating the argument.

  Set $v^l:=E\phi^l$ and $w_j:=Er_j$.  Write
  $\tilde f_j:=e^{ix_j^1\cdot\omega}f_j$ for the re-centred
  sequence, so
  $\tilde u_j:=E\tilde f_j
  =v^1+\sum_{l\geq 2}v^l(\cdot-y_j^l)+w_j$.
  The Euler--Lagrange equation, tested against arbitrary
  $h\in L^2(S^1)$, reads
  \begin{equation}\label{eq:EL-tested}
    \Lambda_j\ip{\tilde f_j}{h}_{L^2(S^1)}
    = \int_{\R^2}|\tilde u_j|^4\tilde u_j\,\overline{Eh}\,dx.
  \end{equation}

  \textbf{Left side.}
  Since $\tilde f_j=\phi^1+\sum_{l\geq 2}e^{-iy_j^l\cdot\omega}\phi^l
  +r_j$ and $\ip{e^{-iy_j^l\cdot\omega}\phi^l}{h}\to 0$ by
  Riemann--Lebesgue ($|y_j^l|\to\infty$), the left side converges to
  $\Lambda_\infty\ip{\phi^1}{h}$.

  \textbf{Right side.}
  Decompose $|\tilde u_j|^4\tilde u_j=|v^1|^4v^1+N_j$, where
  $N_j$ collects all terms containing at least one factor of
  $\eta_j:=\tilde u_j-v^1=\sum_{l\geq 2}v^l(\cdot-y_j^l)+w_j$.
  We show $\int N_j\,\overline{Eh}\to 0$.

  \emph{Step (i): $\eta_j$ vanishes locally in $L^6$.}
  For each fixed $R>0$:
  $\norm{v^l(\cdot-y_j^l)}_{L^6(B(0,R))}
  =\norm{v^l}_{L^6(B(-y_j^l,R))}\to 0$ as $j\to\infty$ (since
  $v^l\in L^6(\R^2)$ and $B(-y_j^l,R)$ escapes to infinity), and
  $\norm{w_j}_{L^6}\to 0$.  Hence
  $\norm{\eta_j}_{L^6(B(0,R))}\to 0$.

  \emph{Step (ii): individual translated-profile terms vanish.}
  Fix $l\geq 2$.  The typical cross term to estimate is
  $I_j^l:=\int|v^1(x)|^4\,v^l(x-y_j^l)\,\overline{Eh(x)}\,dx$.
  Split $\R^2$ into $B(0,R)$ and its complement.
  \begin{itemize}
    \item On $B(0,R)$: $|v^1|\leq(2\pi)^{1/2}$, so
      $|I_j^l|_{B(0,R)}|\leq C\norm{v^l(\cdot-y_j^l)}_{L^6(B(0,R))}
      \norm{Eh}_{L^6}\to 0$ by step~(i).
    \item On $\R^2\setminus B(0,R)$:
      $|I_j^l|_{\R^2\setminus B}|\leq C\norm{v^l}_6
      \norm{Eh}_{L^6(|x|>R)}\to 0$ as $R\to\infty$ (since
      $Eh\in L^6$ has vanishing tails).
  \end{itemize}
  An $\varepsilon/3$ argument in $R$ and $j$ gives $I_j^l\to 0$.

  \emph{Step (iii): the remainder.}
  Terms involving $w_j=Er_j$ are bounded by
  $C\norm{w_j}_6^{}\norm{Eh}_6\to 0$ (using
  $\norm{v^1}_\infty\leq C$ and H\"older).

  \emph{Step (iv): higher-order cross terms.}
  Every monomial in $N_j$ has at least one factor $\eta_j$ or
  $\bar\eta_j$.  The same local-vanishing argument applies: on
  $B(0,R)$, the factor of $\eta_j$ is small in $L^6$; on the
  complement, $Eh$ is small in $L^6$.

  \medskip
  Combining: the right side of~\eqref{eq:EL-tested} converges to
  $\int|v^1|^4v^1\,\overline{Eh}
  =\ip{E^*(|E\phi^1|^4E\phi^1)}{h}$.
  Since $h$ is arbitrary,
  $E^*(|E\phi^1|^4E\phi^1)=\Lambda_\infty\phi^1$.

  The rescaling: if $\alpha_k:=\norm{\phi^k}_2>0$, then
  $\psi^k:=\phi^k/\alpha_k$ satisfies
  $E^*(|E\psi^k|^4E\psi^k)
  =\alpha_k^{-5}E^*(|E\phi^k|^4E\phi^k)
  =\Lambda_\infty\alpha_k^{-5}\phi^k
  =(\Lambda_\infty/\alpha_k^4)\psi^k$.
\end{proof}

\subsection{Existence of a maximiser}

\begin{theorem}[Existence of extremisers for $S^1$ restriction]
  \label{thm:maximiser}
  \proven\quad
  The sharp restriction constant $\mathcal R$ is attained: there exists
  $f_*\in L^2(S^1)$ with $\norm{f_*}_2=1$ and
  $\norm{Ef_*}_{L^6}=\mathcal R$.
\end{theorem}

\begin{proof}
  Let $f_j$ be a maximising sequence: $\norm{f_j}_2=1$,
  $\Lambda_j:=\norm{Ef_j}_6^6\to\mathcal R^6>0$.  Apply the profile
  decomposition (Theorem~\ref{thm:profile-decomp}).

  \textbf{Step 1: at least one profile.}
  Since $\norm{Er_j^K}_6\to 0$ as $K\to\infty$ while
  $\Lambda_j\to\mathcal R^6>0$, the
  decoupling~\eqref{eq:L6-decouple} forces $\phi^1\neq 0$.

  \textbf{Step 2: no splitting.}
  By the restriction inequality, $\norm{E\phi^k}_6^6
  \leq\mathcal R^6\norm{\phi^k}_2^6$.  Set $\alpha_k:=\norm{\phi^k}_2$
  and $a_k:=\alpha_k^2$.  The decoupling gives
  $\mathcal R^6\leq\mathcal R^6\sum_{k}a_k^3$.
  From~\eqref{eq:L2-decouple}, $\sum_k a_k\leq 1$.

  If $M\geq 2$ profiles are nonzero, all $a_k>0$ and $\sum a_k\leq 1$.
  But the multinomial expansion
  $\bigl(\sum_k a_k\bigr)^3
  =\sum_k a_k^3+3\sum_{k\neq l}a_k^2a_l+6\sum_{k<l<m}a_ka_la_m$
  shows $\sum_k a_k^3<\bigl(\sum_k a_k\bigr)^3\leq 1$ (since the
  cross terms are strictly positive).  This gives
  $\mathcal R^6<\mathcal R^6$, a contradiction.

  \textbf{Step 3: convergence.}
  With a single profile: $f_j=g_j^{-1}\cdot\phi^1+r_j$,
  $\norm{r_j}_2\to 0$, $\norm{\phi^1}_2=1$.  Hence
  $g_j\cdot f_j\to\phi^1$ strongly, and $\phi^1$ is a maximiser.
\end{proof}

\begin{remark}[The homogeneity gap]
  The argument uses $6>2$: the $L^6$ norm is superquadratic in~$f$,
  making $\sum a_k^3<(\sum a_k)^3$ when there are multiple profiles.
  This mechanism is specific to maximising sequences and does not
  extend to sub-maximal critical points.
\end{remark}

\subsection{The splitting obstruction for general critical sequences}

For a general sequence of critical points with
$\Lambda_j\to\Lambda_\infty\in[\varepsilon,\mathcal R^6)$, the
strict-convexity argument breaks down: multi-profile decompositions
are \emph{algebraically consistent} with the Euler--Lagrange equation.

\begin{proposition}[Structure of multi-profile decompositions at
  critical sequences]
  \label{prop:multi-profile}
  \proven\quad
  Suppose a sequence of normalised critical points with
  $\Lambda_j\to\Lambda_\infty>0$ admits a profile decomposition
  with $M\geq 2$ nonzero profiles $\phi^1,\ldots,\phi^M$.  Then:
  \begin{enumerate}[leftmargin=2em]
    \item Each normalised profile $\psi^k=\phi^k/\alpha_k$ (where
      $\alpha_k:=\norm{\phi^k}_2$) is a critical point with
      eigenvalue
      $\widetilde\Lambda_k=\Lambda_\infty/\alpha_k^4
      >\Lambda_\infty$.
    \item The number of profiles satisfies
      $M\leq(\mathcal R^6/\varepsilon)^{1/2}$.
    \item The identity $\sum_k\alpha_k^6\widetilde\Lambda_k
      =\Lambda_\infty\sum_k\alpha_k^2=\Lambda_\infty$ holds
      automatically, so the energy decomposition produces no
      contradiction.
  \end{enumerate}
\end{proposition}

\begin{proof}
  (1) Proposition~\ref{prop:profile-critical} gives
  $E^*(|E\phi^k|^4E\phi^k)=\Lambda_\infty\phi^k$.  Rescaling:
  $E^*(|E\psi^k|^4E\psi^k)=(\Lambda_\infty/\alpha_k^4)\psi^k$,
  so $\widetilde\Lambda_k=\Lambda_\infty/\alpha_k^4$.  Since
  $\alpha_k^2<\sum_l\alpha_l^2=1$, we have $\alpha_k<1$ and
  $\widetilde\Lambda_k>\Lambda_\infty$.

  (2) From $\widetilde\Lambda_k\leq\mathcal R^6$:
  $\alpha_k^4\geq\Lambda_\infty/\mathcal R^6
  \geq\varepsilon/\mathcal R^6$, so
  $\alpha_k^2\geq(\varepsilon/\mathcal R^6)^{1/2}$.  The constraint
  $\sum_k\alpha_k^2=1$ gives the bound on~$M$.

  (3) $\sum_k\alpha_k^6\widetilde\Lambda_k
  =\sum_k\alpha_k^6\cdot\Lambda_\infty/\alpha_k^4
  =\Lambda_\infty\sum_k\alpha_k^2=\Lambda_\infty$. \qedhere
\end{proof}

Part~(3) is the central difficulty:
\textbf{the Euler--Lagrange equation alone does not prevent a critical
sequence from splitting into bubbles at diverging locations, each
individually a critical point at a higher eigenvalue level.}
This is the open content of Theorem~\ref{thm:compactness-mod-G}:

\begin{theorem}[Compactness modulo $G$]
  \label{thm:compactness-mod-G}
  \open\quad
  Let $(f_j,\Lambda_j)$ be normalised critical points with
  $\Lambda_j\geq\varepsilon>0$.  Then there exist $g_j\in G$ and a
  subsequence such that $g_j\cdot f_j\to f_\infty$ strongly in
  $L^2(S^1)$ and $\Lambda_j\to\Lambda_\infty=\norm{Ef_\infty}_6^6$.
\end{theorem}

Equivalently: every critical sequence has \textbf{exactly one
nonzero profile} in its decomposition.

\subsection{Approaches to rule out splitting}

We describe three general strategies that may close the gap, and then an
extremiser-specific refinement using nonnegative representatives.

\paragraph{(A) Descending induction on critical values.}
At the maximal level $\Lambda=\mathcal R^6$,
Theorem~\ref{thm:maximiser} proves single-profile concentration.
Suppose inductively that compactness modulo~$G$ holds for critical
sequences with $\Lambda_j\to\Lambda'$ for all
$\Lambda'>\Lambda_\infty$.  A splitting at level~$\Lambda_\infty$
produces profiles at levels $\widetilde\Lambda_k>\Lambda_\infty$,
where the inductive hypothesis already gives finiteness of shapes.
The splitting identity $\sum_k\alpha_k^2=1$,
$\widetilde\Lambda_k=\Lambda_\infty/\alpha_k^4$ imposes algebraic
constraints relating the profile masses to the (now finitely many)
critical values above $\Lambda_\infty$.

\emph{Difficulty}: the base case gives one shape (the maximiser);
but the inductive step requires both finiteness and nondegeneracy at
every intermediate level, creating a dependency between this theorem
and the results of Section~\ref{sec:nondeg}.  Breaking this
circularity requires either a direct nondegeneracy argument at all
levels, or an independent route to single-profile concentration.

\paragraph{(B) Pohozaev--Morawetz monotonicity.}
A localised identity for $u=c(|u|^4u\ast\wh\sigma)$ could yield
a monotonicity formula controlling the spatial distribution of
$|u|^6$ and preventing mass separation.  The starting point would be
to multiply the equation by a radial weight $a(|x|)\,\bar u$ and
integrate; the resulting identity relates $\int a'|u|^6$ to boundary
terms and the nonlocal convolution structure.

\emph{Difficulty}: the convolution against $\wh\sigma$ (which decays
only like $|x|^{-1/2}$) introduces long-range interactions that
resist clean localisation.

\paragraph{(C) Unique continuation for the integrated equation.}
If one can show that two critical points $\phi^1,\phi^2$ at levels
$\widetilde\Lambda_1,\widetilde\Lambda_2>\Lambda_\infty$ cannot
coexist as profiles of a common critical sequence---i.e., the
relation
\[
  f_j\approx e^{-ix_j^1\cdot\omega}\phi^1+e^{-ix_j^2\cdot\omega}\phi^2,
  \qquad |x_j^1-x_j^2|\to\infty,
\]
contradicts $E^*(|Ef_j|^4Ef_j)=\Lambda_jf_j$ at order beyond the
profile-level Euler--Lagrange---this would rule out splitting.  The
relevant obstruction would arise from cross terms at intermediate
spatial scales $|x|\sim|x_j^1-x_j^2|^{1/2}$ where the two Herglotz
waves overlap at polynomial (not exponential) level.

\paragraph{(D) Positive gauge for the extremiser search.}
For the sharp constant, Lemma~\ref{lem:positive-reduction} allows the
search to be restricted to nonnegative critical points.  This is not
available for arbitrary critical sequences, but it removes much of the
translation ambiguity: a nonnegative representative cannot be translated
nontrivially and remain nonnegative.

\begin{lemma}[Nonnegativity fixes the translation parameter]
  \label{lem:positive-gauge}
  Let $0\neq f\in L^2(S^1)$ with $f\geq 0$ a.e.  If
  \[
    e^{i\alpha}e^{-ix_0\cdot\omega}f(R^{-1}\omega)\geq 0
    \qquad\text{for a.e. }\omega\in S^1,
  \]
  then $x_0=0$.  In that case the remaining phase must satisfy
  $e^{i\alpha}=1$ on the support of $f$.
\end{lemma}

\begin{proof}
  The set where $f(R^{-1}\omega)>0$ has positive measure.  On this set,
  the phase $e^{i\alpha}e^{-ix_0\cdot\omega}$ must equal~$1$.  If
  $x_0\neq 0$, the set
  $\{\omega\in S^1:e^{i\alpha}e^{-ix_0\cdot\omega}=1\}$ is finite, since
  it is a finite union of level sets $x_0\cdot\omega=\alpha+2\pi m$.
  This cannot contain a positive-measure subset of~$S^1$.  Hence $x_0=0$.
\end{proof}

The profile-decomposition consequence is the following.  By the
gauge-fixed bootstrap in Proposition~\ref{prop:bootstrap}, a
nonnegative critical sequence with $\Lambda_j\geq\varepsilon$ satisfies
$\norm{f_j}_{C^1(S^1)}\leq C_\varepsilon$.  Thus every escaping profile
vanishes: for $|x_j|\to\infty$ and $h\in C^1(S^1)$, stationary phase gives
\[
  \int_{S^1} e^{ix_j\cdot\omega} f_j(\omega)\overline{h(\omega)}
  \,d\sigma(\omega)
  = O_{C_\varepsilon,h}(|x_j|^{-1/2}),
\]
and density extends this weak convergence to all $h\in L^2(S^1)$.  In the
iterative profile construction, after subtracting finitely many previously
selected profiles, the terms involving
$e^{i(x_j^k-x_j^\ell)\cdot\omega}\phi^\ell$ also vanish by the
Riemann--Lebesgue lemma whenever the translation parameters diverge.
Thus a candidate profile with $|x_j^k|\to\infty$ has zero weak limit.

Thus the extra input needed for the extremiser search is not a new
profile-decomposition mechanism, but the standard gauge-fixed bootstrap:
the $L^2$ normalisation, the equation, and the lower bound on~$\Lambda$
give the uniform norms, while positivity supplies the gauge because pure
translation orbits leave the nonnegative slice by
Lemma~\ref{lem:positive-gauge}.  This closes the escaping-profile
obstruction for the extremiser search without proving compactness of the
full critical set.

\begin{remark}[Sufficiency for identifying $\mathcal R$]
  \label{rem:sufficiency-max}
  For the purpose of \emph{computing} the sharp constant~$\mathcal R$,
  the full Theorem~\ref{thm:compactness-mod-G} is not needed.
  Theorem~\ref{thm:maximiser} gives existence of a maximiser, and if
  the maximiser is unique modulo~$G$ (e.g., the constant function
  $f_0=(2\pi)^{-1/2}$), then $\mathcal R=\norm{Ef_0}_{L^6}$ and
  the problem is solved.  The open compactness theorem would
  additionally yield finiteness of \emph{all} critical values above
  $\varepsilon$, which is a stronger structural result.
\end{remark}

\begin{remark}[Accumulation of critical values at zero]
  \label{rem:accumulation}
  Whether critical values can accumulate at~$0$ is unknown.  The
  compactness + nondegeneracy programme gives finiteness only
  above a fixed threshold.  For the sharp constant, it suffices to
  work above any $\varepsilon$ smaller than the maximum critical
  value.
\end{remark}

% =======================================================================
\section{From the Three Inputs to the Sharp Constant}
\label{sec:finiteness}
% =======================================================================

\subsection{Finiteness after the positivity reduction}

\begin{theorem}[Finiteness of positive critical values in compact windows]
  \label{thm:positive-values-finite}
  \proven\quad
  Fix $0<\Lambda_-<\Lambda_+<\infty$.  Then the set
  \[
    \mathcal V_+(\Lambda_-,\Lambda_+)
    :=
    \left\{
      \Lambda\in[\Lambda_-,\Lambda_+]:
      \begin{array}{l}
      \exists f\in L^2(S^1),\ f\geq 0,\ \norm{f}_2=1,\\
      E^*(|Ef|^4Ef)=\Lambda f
      \end{array}
    \right\}
  \]
  is finite.
\end{theorem}

\begin{proof}
  Let
  \[
    \mathcal C_+
    :=
    \left\{
      (f,\Lambda):
      f\geq0,\ \norm{f}_2=1,\ \Lambda\in[\Lambda_-,\Lambda_+],\
      E^*(|Ef|^4Ef)=\Lambda f
    \right\}.
  \]

  \emph{Step 1: compactness of the positive slice.}
  By the gauge-fixed bootstrap, Proposition~\ref{prop:bootstrap}, for
  every $m\geq0$ there is a constant $C_m=C_m(\Lambda_-)$ such that
  $\norm{f}_{C^m(S^1)}\leq C_m$ for every $(f,\Lambda)\in\mathcal C_+$.
  The reason this applies without an additional translation choice is
  precisely Lemma~\ref{lem:positive-gauge}: the nonnegative slice fixes
  the noncompact translation parameter.

  Given a sequence $(f_j,\Lambda_j)\in\mathcal C_+$, pass to a subsequence
  with $\Lambda_j\to\Lambda_\infty\in[\Lambda_-,\Lambda_+]$.  The uniform
  $C^{m+1}$ bound and Arzel\`a--Ascoli give, after passing to a further
  subsequence, $f_j\to f_\infty$ in $C^m(S^1)$ for each fixed~$m$ (and in
  particular in $L^2$).  Then $f_\infty\geq0$ and
  $\norm{f_\infty}_2=1$.  Since
  \[
    T(f):=E^*(|Ef|^4Ef)
  \]
  is continuous $L^2(S^1)\to L^2(S^1)$ by Stein--Tomas and H\"older,
  passing to the limit in $T(f_j)=\Lambda_j f_j$ gives
  $T(f_\infty)=\Lambda_\infty f_\infty$.  Hence
  $(f_\infty,\Lambda_\infty)\in\mathcal C_+$.  Thus $\mathcal C_+$ is
  compact in $L^2(S^1)\times\R$.

  \emph{Step 2: local finite-dimensional analytic structure.}
  Work on the real Hilbert space $H:=L^2(S^1;\C)$ with real inner product
  $\re\ip{\cdot}{\cdot}$.  Define
  \[
    \Phi:H\times\R\to H\times\R,\qquad
    \Phi(f,\Lambda):=\bigl(T(f)-\Lambda f,\ \norm{f}_2^2-1\bigr).
  \]
  The map $T$ is a real-analytic quintic polynomial map:
  $T(f)=E^*((Ef)^3(\overline{Ef})^2)$.  At a zero $(f_0,\Lambda_0)$,
  \[
    D\Phi_{(f_0,\Lambda_0)}(h,\mu)
    =
    \left(
      (K_{u_0}-\Lambda_0 I)h-\mu f_0,\
      2\re\ip{f_0}{h}
    \right),
    \qquad u_0=Ef_0.
  \]
  By Corollary~\ref{cor:fredholm}, $K_{u_0}-\Lambda_0 I$ is Fredholm of
  index~$0$.  The displayed derivative is a finite-rank perturbation of
  $(h,\mu)\mapsto((K_{u_0}-\Lambda_0 I)h,\mu)$, and is therefore Fredholm
  of index~$0$.

  The analytic Fredholm Lyapunov--Schmidt reduction now applies: near
  each zero of~$\Phi$, the zero set $\Phi^{-1}(0)$ is homeomorphic to the
  zero set of a finite-dimensional real-analytic map.  By the standard
  local structure theorem for real-analytic sets (equivalently,
  Lojasiewicz local finiteness), this local zero set has only finitely
  many connected components, and each such component is locally
  arc-connected through finitely many analytic arcs.

  \emph{Step 3: the value is constant on local critical components.}
  Let $J(f):=\norm{Ef}_{L^6}^6$.  On $\Phi^{-1}(0)$ one has
  $\Lambda=J(f)$, because
  \[
    \Lambda\norm{f}_2^2
    =\re\ip{T(f)}{f}
    =\int_{\R^2}|Ef(x)|^6\,dx.
  \]
  If $t\mapsto(f(t),\Lambda(t))$ is a $C^1$ arc in $\Phi^{-1}(0)$, then
  $\norm{f(t)}_2=1$ and $T(f(t))=\Lambda(t)f(t)$, so
  \[
    \frac{d}{dt}J(f(t))
    =6\re\ip{T(f(t))}{f'(t)}
    =6\Lambda(t)\re\ip{f(t)}{f'(t)}
    =0.
  \]
  Therefore $J(f(t))$, and hence $\Lambda(t)$, is constant on every such
  arc.  Since the local analytic components are connected by finite
  chains of these arcs, $\Lambda$ is constant on each local connected
  component of $\Phi^{-1}(0)$.

  Finally, compactness of $\mathcal C_+$ gives a finite cover by
  neighbourhoods of the type described in Step~2.  Each neighbourhood
  contains only finitely many local critical components, and Step~3 gives
  at most one value of~$\Lambda$ on each of them.  Hence the projection
  of $\mathcal C_+$ onto the $\Lambda$ coordinate, namely
  $\mathcal V_+(\Lambda_-,\Lambda_+)$, is finite.
\end{proof}

\subsection{A spectral gap around the constant value}

Let
\[
  f_0=(2\pi)^{-1/2},\qquad
  \Lambda_0:=\norm{Ef_0}_{L^6}^6
  =16\pi^4\int_0^\infty J_0(r)^6\,r\,dr.
\]
If $C_{\mathrm{ST}}$ denotes any explicit Stein--Tomas bound for
$E:L^2(S^1)\to L^6(\R^2)$, then every normalised critical value satisfies
$\Lambda\leq C_{\mathrm{ST}}^6$.  Thus, when working with the eigenvalue
parameter~$\Lambda$ rather than with the restriction norm
$\Lambda^{1/6}$, the natural finite window for the extremiser search is
\[
  W_0:=\left[\frac12\Lambda_0,\ C_{\mathrm{ST}}^6\right].
\]
By Theorem~\ref{thm:positive-values-finite},
$\mathcal V_+(W_0):=\mathcal V_+(\Lambda_0/2,C_{\mathrm{ST}}^6)$ is a
finite set and contains~$\Lambda_0$.  Hence the nonlinear spectral gap
around the constant value,
\begin{equation}\label{eq:positive-gap-def}
  \gamma_0
  :=
  \operatorname{dist}\bigl(
    \Lambda_0,\,
    \mathcal V_+(W_0)\setminus\{\Lambda_0\}
  \bigr),
\end{equation}
is strictly positive, with the convention $\gamma_0=+\infty$ if
$\mathcal V_+(W_0)=\{\Lambda_0\}$.

The proof of Theorem~\ref{thm:positive-values-finite} gives positivity of
$\gamma_0$, but not a numerical lower bound.  To make the gap quantitative,
one needs to quantify two pieces.

\paragraph{Local isolation at the constant.}
On the slice orthogonal to phase and translations, define the linear
nondegeneracy gap
\[
  \nu_0:=
  \min\left\{
    \frac45\Lambda_0,\;
    \inf_{n\geq2}\abs{\mathcal I_n-\Lambda_0},\;
    \inf_{n\geq2}\abs{5\mathcal I_n-\Lambda_0}
  \right\},
  \qquad
  \mathcal I_n=16\pi^4\mathcal J_n.
\]
If interval arithmetic proves $\nu_0\geq\underline\nu>0$ and
$M_0(r)$ is a Lipschitz constant for the derivative of the sliced
Euler--Lagrange map in the $L^2$ ball
$\norm{f-f_0}_2\leq r$, then the Newton--Kantorovich estimate
\[
  r_0\leq \frac{\underline\nu}{2M_0(r_0)}
\]
certifies that the constant is the only positive critical point in that
ball, up to the trivial rotation symmetry of the constant.

A rough numerical quadrature of the Bessel integrals gives
\[
  \int_0^\infty J_0(r)^6r\,dr \approx 0.33678766,\qquad
  \Lambda_0\approx 524.899,\qquad
  \Lambda_0^{1/6}\approx 2.84018,
\]
and suggests
\[
  \frac{\nu_0}{\Lambda_0}
  =
  \min_{n\geq2}
  \left\{
    \abs{\frac{\mathcal I_n}{\Lambda_0}-1},
    \abs{\frac{5\mathcal I_n}{\Lambda_0}-1}
  \right\}
  \approx
  0.452,
\]
with the minimum apparently coming from the $5\mathcal I_2-\Lambda_0$
branch.  A reasonable first rigorous target is therefore
\begin{equation}\label{eq:nu-target}
  \nu_0\geq 0.4\,\Lambda_0.
\end{equation}
This target only requires a finite interval-arithmetic check for the
small modes, together with the already proved tail estimate
$\mathcal J_n\to0$ for large~$n$.

\paragraph{Current certified finite-mode check.}
The script \texttt{numerics/circle\_gap\_flint\_verify.py} gives a first
ball-arithmetic certificate for the small-mode part of this target.  It uses
\texttt{python-flint} Arb balls, evaluates the Bessel functions at the
midpoints of the radial mesh, and encloses the variation on each interval by
the elementary bounds $\abs{J_m}\leq1$ and $\abs{J'_m}\leq1$ for integer
$m$ on the real line.  With cutoff $R=200$ and the conditional tail majorant
$\abs{J_m(r)}\leq (2/(\pi r))^{1/2}$ on $[R,\infty)$, the current runs prove
\[
  \frac{\mathcal J_2}{\mathcal J_0}\leq 0.117364,\qquad
  \frac{\mathcal J_3}{\mathcal J_0}\leq 0.091870,\qquad
  \frac{\mathcal J_n}{\mathcal J_0}\leq 0.072366
  \quad (4\leq n\leq12).
\]
The first number is the bottleneck and was obtained with mesh
$h=0.0025$; the other displayed bounds use $h=0.01$.  In particular, for
$2\leq n\leq12$ the finite-mode certificate proves
\[
  \frac{\mathcal J_n}{\mathcal J_0}<0.12,
\]
and hence
\[
  \abs{5\mathcal I_n-\Lambda_0}>0.4\Lambda_0,\qquad
  \abs{\mathcal I_n-\Lambda_0}>0.4\Lambda_0
  \qquad (2\leq n\leq12).
\]
What remains for a complete proof of~\eqref{eq:nu-target} is to replace or
quote the tail majorant in a final form and prove the large-mode estimate
for all $n\geq13$.

\paragraph{Global separation from the rest of the positive window.}
After the local ball $\norm{f-f_0}_2<r_0$ has been certified, the remaining
set
\[
  \mathcal C_+^{\mathrm{far}}
  :=
  \left\{
    (f,\Lambda):
    \begin{array}{l}
      f\geq0,\ \norm{f}_2=1,\ E^*(|Ef|^4Ef)=\Lambda f,\\
      \Lambda\in W_0,\ \norm{f-f_0}_2\geq r_0
    \end{array}
  \right\}
\]
is compact.  A validated search on this compact set should return finitely
many disjoint value intervals $B_1,\ldots,B_N$ containing every positive
critical value in the far region.  Then the explicit certificate
\[
  \gamma_{\mathrm{cert}}
  :=
  \min_{1\leq j\leq N,\ \Lambda_0\notin B_j}
  \operatorname{dist}(\Lambda_0,B_j)
\]
is a rigorous lower bound for~$\gamma_0$.

Thus the current proof gives the qualitative gap~\eqref{eq:positive-gap-def}.
The practical quantitative route is: certify the local linear
target~\eqref{eq:nu-target}, use it to remove a ball around the constant
by Newton--Kantorovich, and then use the interval search on the compact
far region to produce $\gamma_{\mathrm{cert}}$.

\subsection{What a rigorous brute-force search would mean}
\label{sec:bruteforce-search}

The phrase ``search the compact region'' should not be interpreted as
requiring a prior description of the compact set.  The algorithm searches
a computable \emph{outer enclosure} of that compact set.  The compactness
theorem supplies the analytic reason such an enclosure can be made finite;
the actual numerical proof must supply explicit constants.

There are two complementary brute-force targets.

\paragraph{Eigenvalue exclusion.}
Fix an interval $I\subset W_0$ for~$\Lambda$.  A rigorous search attempts
to prove:
\[
  \nexists f\geq0,\quad \norm{f}_2=1,\quad
  E^*(|Ef|^4Ef)=\Lambda f,\quad \Lambda\in I,
\]
except possibly inside boxes already validated as containing known
solutions.  The output of such a search is not ``success/failure'' in the
usual numerical sense.  Each box in the finite cover is labelled
\emph{empty}, \emph{validated}, or \emph{undecided}.  A theorem is obtained
only if no undecided boxes remain.

\paragraph{Direct upper-bound search.}
Instead of solving the Euler--Lagrange equation, one can try to prove the
stronger variational statement
\[
  \sup_{\substack{f\geq0\\ \norm{f}_2=1}}
  \norm{Ef}_{L^6}^6
  \leq \Lambda_0+\eta.
\]
If an eigenvalue gap $\gamma_0>\eta$ has already been certified in the
positive window, then this upper bound forces the maximiser to have value
$\Lambda_0$.  This route may be computationally easier because it asks for
upper bounds on a scalar polynomial functional, rather than for a complete
classification of roots of a vector equation.

\paragraph{Fourier truncation and tails.}
Write
\[
  f(\theta)=\sum_{n\in\Z}a_ne^{in\theta},\qquad
  P_Nf=\sum_{|n|\leq N}a_ne^{in\theta},\qquad r_N=f-P_Nf.
\]
For real-valued nonnegative candidates, $a_{-n}=\overline{a_n}$; if one
also imposes an even symmetry ansatz, the unknowns reduce further to real
cosine coefficients.  The bootstrap gives an explicit tail estimate of
the form
\[
  \norm{r_N}_{L^2}\leq \tau_N(B_m),
  \qquad
  \tau_N(B_m)\to0,
\]
where $B_m$ is a certified $C^m$ or Sobolev bound for positive critical
points in the window.  Thus one searches over the finite coefficient vector
$(a_{-N},\ldots,a_N)$ and an interval uncertainty representing~$r_N$.

The finite-dimensional residual is
\[
  F_k(a,\Lambda)
  :=
  \widehat{T(P_Nf)}(k)-\Lambda a_k,
  \qquad |k|\leq N,
\]
with an explicit interval error for all terms containing at least one
factor of~$r_N$.  A convenient crude tail bound follows from
Stein--Tomas and the Lipschitz estimate for $z\mapsto |z|^4z$:
\[
  \norm{T(f)-T(P_Nf)}_2
  \leq
  C\,C_{\mathrm{ST}}^6
  \bigl(\norm{f}_2+\norm{P_Nf}_2\bigr)^4
  \norm{r_N}_2,
\]
with an absolute constant~$C$.  This is probably not sharp enough for the
final proof, but it is the correct first enclosure.  The Bessel-product
technology should be used later to reduce this overestimate.

\paragraph{Why sixfold Bessel products appear.}
The extension of one Fourier mode is
\[
  E(e^{in\theta})(r,\phi)
  =2\pi i^{|n|}J_{|n|}(r)e^{in\phi}.
\]
Consequently, the coefficient $\widehat{T(P_Nf)}(k)$ is a finite quintic
sum over indices satisfying the angular selection rule
\[
  n_1+n_2+n_3-n_4-n_5-k=0,
\]
with constants given by radial integrals of six Bessel functions:
\[
  \int_0^\infty
  J_{|n_1|}(r)J_{|n_2|}(r)J_{|n_3|}(r)
  J_{|n_4|}(r)J_{|n_5|}(r)J_{|k|}(r)\,r\,dr.
\]
These constants must be enclosed rigorously.  This is exactly the type of
input studied by Oliveira e Silva--Thiele in their estimates for sixfold
Bessel products, and by the band-limited maximiser papers described below.

\paragraph{Interval branch-and-bound.}
The finite search region is a box in
\[
  (a_{-N},\ldots,a_N,\Lambda),
\]
augmented by interval constraints for normalisation, positivity, and the
tail.  Positivity can be enforced by interval evaluation on arcs: if
$P_Nf$ is bounded below by more than the certified tail and mesh
interpolation error, the box is positive; if it is forced negative
somewhere, the box is discarded; otherwise it is subdivided.

For the eigenvalue search, each box is processed as follows.
\begin{enumerate}[leftmargin=2em]
  \item Evaluate the interval residuals $F_k$ and the normalisation
    residual.  If some residual interval excludes~$0$, discard the box.
  \item Apply an interval Newton or Krawczyk operator to the finite system.
    If it proves no zero lies in the box, discard it.  If it maps strictly
    into the box, validate a unique zero.
  \item If neither test succeeds, subdivide the box, unless the box is
    already below the target resolution; in that case it is marked
    undecided and no theorem is claimed.
\end{enumerate}
For the variational upper-bound search, the analogous step is interval
upper-bounding of the scalar functional $\norm{E(P_Nf)}_6^6$ plus the
tail error.  Boxes whose upper bound is below $\Lambda_0+\eta$ are
discarded; boxes near the top are subdivided or handled by local concavity
at the constant.

\paragraph{How this could prove the sharp constant.}
A plausible staged proof would be:
\begin{enumerate}[leftmargin=2em]
  \item Certify the local linear gap target
    $\nu_0\geq0.4\Lambda_0$ by interval Bessel integration.
  \item Use Newton--Kantorovich to remove an explicit ball around~$f_0$.
  \item Run the eigenvalue exclusion search on
    $[\Lambda_0,\Lambda_0+\gamma]$ outside that ball.  If it exhausts the
    boxes, then no positive critical value lies in this interval.
  \item Separately prove by branch-and-bound that
    $\norm{Ef}_6^6\leq\Lambda_0+\gamma/2$ for all positive normalised
    candidates in the compact window.
  \item Since an extremiser exists and may be chosen nonnegative, the
    maximising value must then equal~$\Lambda_0$.
\end{enumerate}
The same architecture can be run in restriction-norm units
$\Lambda^{1/6}$.  One should be careful about scales: near~$\Lambda_0$,
a gap of $10^{-3}$ in $\Lambda^{1/6}$ corresponds to a gap of order
$6\Lambda_0^{5/6}\cdot10^{-3}$ in~$\Lambda$.

\paragraph{Relevant numerical and analytic references.}
The closest existing works are:
\begin{itemize}[leftmargin=2em]
  \item D. Oliveira e Silva and C. Thiele,
    \emph{Estimates for certain integrals of products of six Bessel
    functions}, Rev. Mat. Iberoam. 33 (2017), no. 4, 1423--1462,
    \href{https://doi.org/10.4171/RMI/977}{doi:10.4171/RMI/977}.
    This is the main reference for rigorous estimates of the sixfold
    Bessel integrals needed in the residual coefficients.
  \item E. Carneiro, D. Foschi, D. Oliveira e Silva and C. Thiele,
    \emph{A sharp trilinear inequality related to Fourier restriction on
    the circle}, Rev. Mat. Iberoam. 33 (2017), no. 4, 1463--1486,
    \href{https://doi.org/10.4171/RMI/978}{doi:10.4171/RMI/978}.
    This uses the sixfold Bessel estimates and proves, among other things,
    that constants are local extremizers for the circle Tomas--Stein
    problem.
  \item D. Oliveira e Silva, C. Thiele and P. Zorin-Kranich,
    \emph{Band-limited maximizers for a Fourier extension inequality on
    the circle}, Exp. Math. 31 (2022), no. 1, 192--198,
    \href{https://doi.org/10.1080/10586458.2019.1596847}
    {doi:10.1080/10586458.2019.1596847}.  They prove uniqueness of
    constants among real-valued functions with Fourier modes up to
    degree~$30$.
  \item J. Barker, C. Thiele and P. Zorin-Kranich,
    \emph{Band-limited maximizers for a Fourier extension inequality on
    the circle, II}, Exp. Math. 32 (2023), no. 2, 280--293,
    \href{https://doi.org/10.1080/10586458.2021.1926011}
    {doi:10.1080/10586458.2021.1926011}.  This pushes the band-limited
    computation to modes up to degree~$120$.
  \item R. E. Moore, R. B. Kearfott and M. J. Cloud,
    \emph{Introduction to Interval Analysis}, SIAM, 2009.  This is a
    standard reference for interval Newton and Krawczyk methods.
  \item F. Johansson, \emph{Arb: efficient arbitrary-precision
    midpoint-radius interval arithmetic}, IEEE Trans. Comput. 66 (2017),
    no. 8, 1281--1292,
    \href{https://doi.org/10.1109/TC.2017.2690633}
    {doi:10.1109/TC.2017.2690633}.  Arb/ball arithmetic is a natural
    implementation platform for rigorous Bessel-integral and polynomial
    residual enclosures.
\end{itemize}

\begin{theorem}[Conditional finiteness of shapes]
  \label{thm:finiteness}
  \conditional\quad
  Assume:
  \begin{enumerate}[leftmargin=2em]
    \item Compactness modulo $G$ (Theorem~\ref{thm:compactness-mod-G})
      holds for critical points with $\Lambda\geq\varepsilon$;
    \item Every such critical point is nondegenerate modulo $G$
      (i.e., $\ker L=T_{f_0}(G\cdot f_0)$ on the tangent space to the
      unit sphere).
  \end{enumerate}
  Then the set of shapes $\{[f,\Lambda]\in\mathcal Z:\Lambda\geq\varepsilon\}$
  is finite.  In particular, the set of critical values above $\varepsilon$
  is finite.
\end{theorem}

\begin{proof}
  By compactness modulo~$G$, the quotient
  $\mathcal Z_\varepsilon:=\{f:\text{critical},\,\Lambda\geq\varepsilon\}/G$
  is compact.

  By nondegeneracy: at each $[f_0]\in\mathcal Z_\varepsilon$, choose a slice
  $S_{f_0}$ complementary to $T_{f_0}(G\cdot f_0)$ in the tangent space
  to $\norm{f}_2=1$.  The condition $\ker L|_{S_{f_0}}=\{0\}$ plus
  Fredholm index zero implies $L|_{S_{f_0}}$ is invertible with bounded
  inverse $\norm{L^{-1}|_{S_{f_0}}}\leq\mu_{f_0}^{-1}$.  By the
  implicit function theorem, $[f_0]$ is an isolated point of
  $\mathcal Z_\varepsilon$.

  A compact set all of whose points are isolated is finite.
\end{proof}

\begin{corollary}[Identification of the sharp constant]
  Under the hypotheses of Theorem~\ref{thm:finiteness}, the sharp
  restriction constant is
  \[
    \mathcal R = \max\{\Lambda_j^{1/6}:\Lambda_j\geq\varepsilon\}
  \]
  for any $\varepsilon$ smaller than all critical values in the finite list.
  The maximiser is the shape achieving the largest~$\Lambda_j$.
\end{corollary}

\subsection{What remains for the sharp constant on $S^1$}

Two paths are available, depending on the desired strength of the result.

\paragraph{Path 1: Identify $\mathcal R$ (the sharp constant).}
\begin{enumerate}[leftmargin=2em]
  \item \textbf{Verify Bessel integrals.}  Check
    $\mathcal J_n\neq\mathcal J_0$ and $5\mathcal J_n\neq\mathcal J_0$
    for $n=2,3,\ldots,n_*$ by interval arithmetic.  This establishes
    nondegeneracy at the constant solution.  \emph{Routine numerical
    computation.}

  \item \textbf{Restrict the search to nonnegative eigenfunctions.}
    By Lemma~\ref{lem:positive-reduction}, any extremiser can be replaced
    by a nonnegative extremiser, possibly in a different shape.  Therefore,
    for the sharp constant, one only needs to retain those candidate
    values~$\Lambda_j$ for which the nonlinear eigenvalue equation has at
    least one nonnegative solution.  The positive gauge also removes exact
    translation ambiguity by Lemma~\ref{lem:positive-gauge}, and
    Theorem~\ref{thm:positive-values-finite} proves that the retained
    critical values in any fixed positive window are finite.

  \item \textbf{Show the constant is the unique maximiser.}
    Theorem~\ref{thm:maximiser} gives existence of a maximiser.
    If one can prove that the constant function $f_0=(2\pi)^{-1/2}$
    is the unique maximiser modulo~$G$---e.g., by a rearrangement
    or symmetrisation argument, or by ruling out non-radial
    maximisers via the eigenvalue analysis---then
    $\mathcal R=\norm{Ef_0}_{L^6}$ and the programme is complete.
    \emph{This bypasses compactness of the full critical set.}
\end{enumerate}

\paragraph{Path 2: Finiteness of all critical values above
$\varepsilon$ (the full structural result).}
\begin{enumerate}[leftmargin=2em]
  \item Bessel integrals (same as above).

  \item \textbf{Rule out multi-profile splitting
    (Theorem~\ref{thm:compactness-mod-G}).}
    The profile decomposition is in place
    (Theorem~\ref{thm:profile-decomp}), and each profile inherits the
    Euler--Lagrange equation
    (Proposition~\ref{prop:profile-critical}).  The open step is to
    show that no critical sequence admits more than one profile.
    Section~\ref{sec:compactness} identifies three general strategies:
    descending induction, Pohozaev--Morawetz monotonicity, and unique
    continuation for cross terms.  \emph{Requires new analytic input.}

  \item \textbf{Enumerate non-constant shapes (if any).}  Either prove
    that the constant is the only critical shape, or verify
    nondegeneracy at all other shapes.
\end{enumerate}

% =======================================================================
% PART III: EXTENSIONS AND SUPPLEMENTARY MATERIAL
% =======================================================================

\part{Extensions and Supplementary Material}

% =======================================================================
\section{The Higher-Dimensional Transversality Route}
\label{sec:higher-dim}
% =======================================================================

For $n\geq 2$ and the equation $\rho=cT(\rho)$ on $L^2(S^{n-1})$,
Proposition~\ref{prop:support-forcing} forces $\wh u$ onto $S^{n-1}$,
and the problem reduces to the incidence-fibre equation~\eqref{eq:sphere-eq}.
The symmetry group is $G=\R^n\rtimes SO(n)$, acting on densities by
$(x_0,R)\cdot\rho(\omega)=e^{-ix_0\cdot\omega}\rho(R^{-1}\omega)$.

\subsection{The five-step programme}

The route to finiteness of shapes requires five analytic inputs.  We state
each and assess its status for general~$n$.

\paragraph{Step 1: Smoothing / Fredholm estimate.}  \open\quad
The target is a multilinear smoothing estimate
\[
  A_k:H^s(S^{n-1})^k\to H^{s+\varepsilon}(S^{n-1}),
  \qquad
  \norm{A_k(\rho_1,\ldots,\rho_k)}_{H^{s+\varepsilon}}
  \leq C_s\prod_{j=1}^k\norm{\rho_j}_{H^s},
\]
for some $\varepsilon>0$, where $A_k$ is the $k$-linear operator behind
the incidence-fibre equation.  A weaker but sufficient target is
compactness of $A_k$ on bounded sets.

For the $S^1$ problem, the truncation argument
(Theorem~\ref{thm:K-compact}) bypasses this step entirely.  For $n\geq 2$,
the proof must exploit the curvature of $S^{n-1}$ and the transversality
of the constraint $\omega_1+\cdots+\omega_k=\omega$.  The key obstruction
is the existence of \emph{degenerate fibres}; see
Section~\ref{sec:degenerate-fibres}.

\paragraph{Step 2: Regularity bootstrap.}  \open\quad
Show that weak $L^2$ solutions are automatically smooth:
$\rho\in L^2(S^{n-1}),\;\rho=cT(\rho)\implies\rho\in H^s$ for all~$s$.
If Step~1 gives $T:L^2\to H^\varepsilon$, iteration yields the bootstrap.

\paragraph{Step 3: Local slices for the symmetry action.}  \open\quad
At a solution $\rho$, construct a complement $S_\rho$ to the orbit tangent
\[
  \mathfrak g\cdot\rho
  = \{-i(a\cdot\omega)\rho-(A\omega)\cdot\nabla_{S^{n-1}}\rho:
  a\in\R^n,\;A\in\mathfrak{so}(n)\}.
\]
This must be done separately on orbit-type strata (radial, zonal, fully
asymmetric solutions have different stabiliser dimensions).

\paragraph{Step 4: Nondegeneracy modulo symmetry.}  \open\quad
Verify $\ker L_\rho=\mathfrak g\cdot\rho$ at each solution, where
$L_\rho=I-ckA_k(\rho,\ldots,\rho,\cdot)$.  For the $S^1$ problem at
the constant solution, this reduces to Bessel integrals
(Theorem~\ref{thm:nondeg-criterion}).

\paragraph{Step 5: Compactness modulo~$G$.}  \open\quad
Prove sequential compactness of the solution set modulo~$G$:
\[
  \rho_j=cT(\rho_j)\implies\exists g_j\in G:\;g_j\cdot\rho_j\to\rho_\infty.
\]
This requires a priori bounds ($\norm\rho\leq R$), a non-collapse lower
bound ($\norm\rho\geq r>0$), and a no-escape condition for translations.
Rotations are compact; translations are the genuine noncompact obstruction.

\subsection{Degenerate fibres}
\label{sec:degenerate-fibres}

The fibre $V_k(\omega)=\{(\omega_1,\ldots,\omega_k)\in(S^{n-1})^k:
\sum\omega_j=\omega\}$ is generically a smooth manifold of dimension
$(n-1)k-(n-1+1)=(n-1)(k-1)-1$... when the differential of
$\Sigma:\omega_1+\cdots+\omega_k$ has full rank.  Degeneracies occur when
all tangent lines $T_{\omega_j}S^{n-1}$ are parallel.

\textbf{Example for $k=5$ on $S^1$:}  The configuration
$(\omega_0,\omega_0,\omega_0,-\omega_0,-\omega_0)$ with output
$\omega=\omega_0$ is degenerate.  Here
$d\Sigma=3(\text{tangent at }\omega_0)-2(\text{tangent at }-\omega_0)$
has rank~$1$ instead of~$2$ (the tangents at $\omega_0$ and $-\omega_0$
are parallel on $S^1$).  This is a codimension-$1$ degeneracy, producing
an integrable singularity $\sim|t|^{-1/2}$ in the fibre density.

In higher dimensions, the degenerate locus can have higher codimension
(tangents at $\omega_0$ and $-\omega_0$ are parallel in $\R^n$ only
when $n=1$; for $n\geq 2$, anti-parallel tangent \emph{spaces} are
generically transverse).  This suggests that the smoothing estimate may
be easier for $n\geq 3$ than for $n=2$.

\subsection{Orbit dimension and radial rigidity}

\begin{proposition}[Orbit dimension, $n=3$]
  For $G=\R^3\rtimes SO(3)$ ($\dim G=6$), the orbit dimension of a
  solution~$u$ is $6-\dim\Stab_G(u)$.  Specifically:
  \begin{itemize}
    \item Radial: $d=3$ (only translations).
    \item Zonal (one axis of symmetry): $d=5$.
    \item Fully asymmetric: $d=6$.
  \end{itemize}
\end{proposition}

\begin{conjecture}[Radial rigidity]
  \label{conj:radial}
  For the equation $u=c(|u|^{2(k-1)}u\ast\wh\sigma)$ on $\R^3$, every
  non-trivial solution is rotationally invariant about some point (its
  ``centre of mass'').  If true, all solutions belong to $3$-dimensional
  translation orbits, and the shape count reduces to counting radial profiles.
\end{conjecture}

\begin{warning}
  Conjecture~\ref{conj:radial} is currently \emph{unproved}.  The
  preceding sections on the $n=3$ problem treated it as established; it
  should be regarded as a useful heuristic until proved.
\end{warning}

% =======================================================================
\section{Supplementary: Gegenbauer Expansion and Galerkin Models}
\label{sec:galerkin}
% =======================================================================

\subsection{Gegenbauer expansion}
\heuristic

For completeness, we record the formal harmonic decomposition.
Let $\rho(\omega)=\sum_\ell a_\ell C_\ell^\lambda(\omega\cdot\eta)$ with
$\lambda=(n-2)/2$.  The product of Gegenbauer polynomials linearises as
$C_{\ell_1}^\lambda C_{\ell_2}^\lambda
=\sum_j\gamma_j(\ell_1,\ell_2)C_j^\lambda$.
If $T$ were a spherical convolution, the Funk--Hecke theorem would give
$C_\ell^\lambda\ast_S\sigma=\lambda_\ell(\sigma)C_\ell^\lambda$.
Substituting into $\rho=cT(\rho)$ and equating coefficients would give an
algebraic system for $\{a_\ell\}$.

\textbf{However}, as noted in Warning~\ref{warn:not-convolution}, the
incidence-fibre operator is \emph{not} spherical convolution.  The
Gegenbauer expansion gives a useful \emph{model} but does not rigorously
reduce to finitely many coefficients without separate justification.

\subsection{Galerkin approximation}

Let $H_L=\operatorname{span}\{Y_{\ell m}:0\leq\ell\leq L\}$ and $P_L$ the
$L^2(S^{n-1})$ projection.  The projected equation $P_LF(\rho_L)=0$ is a
polynomial system of degree~$k$ in $N(L)=(L+1)^2$ unknowns (for $n=3$).

\begin{proposition}[Galerkin shape count]
  After imposing algebraic slice conditions, the number of Galerkin shapes
  at truncation~$L$ is bounded by the BKK mixed volume of the sliced
  Newton polytopes, in particular by $k^{N(L)}$ (B\'ezout).
\end{proposition}

\begin{warning}[Limit caveat]
  This is a statement about the fixed-$L$ model.  It does not imply
  finiteness of true shapes as $L\to\infty$ without uniform a priori bounds,
  no-escape compactness, and stability of roots under high-mode
  perturbation.
\end{warning}

\begin{remark}[How Galerkin helps]
  A Galerkin computation becomes rigorous if paired with a high-mode bound
  $\norm{(I-P_L)T(\rho)}_X\leq\varepsilon_L\Phi(\norm\rho_X)$ with
  $\varepsilon_L\to 0$, plus an inverse bound on the sliced Jacobian at each
  computed root.  A Newton--Kantorovich argument would then validate nearby
  true solutions and certify completeness.
\end{remark}

% =======================================================================
\section{Supplementary: Validated Search Strategy}
\label{sec:validated-search}
% =======================================================================

Assuming the transversality route is completed (modulo~$G$ the solution set
is finite), the next goal is to \emph{find every shape} by a finite,
checkable search.

\subsection{Constants needed}

\begin{enumerate}[leftmargin=2em]
  \item \textbf{A priori radius:} $\norm\rho_X\leq R(c,k)$.
  \item \textbf{Non-collapse:} $\norm\rho_X\geq r(c,k)>0$.
  \item \textbf{Regularity/mesh modulus:} $\norm\rho_{C^m}\leq B_m$.
  \item \textbf{Validated quadrature:} error $q_h(B_m)\to 0$ for the
    fibre integral on a mesh of size~$h$.
  \item \textbf{Slice conditioning:}
    $\norm{L_\rho\eta}_X\geq\mu(\rho)\norm\eta_X$ on $S_\rho$, with
    $\mu_*:=\inf_\rho\mu(\rho)>0$.
  \item \textbf{Lipschitz:}
    $\norm{DF(\rho)-DF(\tilde\rho)}\leq M\norm{\rho-\tilde\rho}_X$.
\end{enumerate}

\subsection{Separation and packing}

From the conditioning bound, the quantitative isolation scale is
$\delta_*\simeq\mu_*/M$: two distinct shapes cannot lie within
$\delta_*$ of each other in the slice.  A crude packing bound is
\[
  \#\{\text{shapes}\}
  \leq
  \left(1+\frac{2R_h}{\delta_*-2\eta_h-2q_h}\right)^{N_{\mathrm{mesh}}(h)},
\]
provided $\delta_*>2\eta_h+2q_h$.

\subsection{The search pipeline}

\begin{enumerate}[leftmargin=2em]
  \item Choose a quasi-uniform mesh on $S^{n-1}$ (or directly on $S^1$
    for the circle problem) and fix gauge conditions.
  \item Represent candidates by interval boxes for nodal values plus global
    regularity bounds.
  \item For each box, enclose the residual $F(\rho)$ at mesh points using
    validated quadrature.
  \item Discard boxes with nonzero residual enclosure.
  \item Apply Krawczyk/interval Newton to surviving boxes.
  \item Validate uniqueness and lift to the full space.
  \item Certify completeness: every box discarded or assigned.
\end{enumerate}

\subsection{Main difficulties}

\begin{enumerate}[leftmargin=2em]
  \item \textbf{Translation gauge:} $e^{-ix\cdot\omega}\rho(\omega)$ creates
    oscillation as $|x|\to\infty$.  A canonical centre must be fixed.
  \item \textbf{Uniform conditioning:} $\mu_*>0$ may be the hardest
    constant to obtain without knowing all solutions.
  \item \textbf{Mesh consistency:} nodal errors must not hide solutions
    between mesh points.
  \item \textbf{Dimension:} the effective search dimension grows with
    mesh size.  Symmetry decomposition and adaptive refinement are essential.
\end{enumerate}

\subsection{Specialisation to $S^1$}

For the circle problem, the validated search simplifies considerably:
\begin{itemize}
  \item The mesh is a grid on $[0,2\pi)$ with $N$ equally spaced points.
  \item The gauge fixes phase, rotation of $S^1$, and one translation.
  \item The fibre integral reduces to a trigonometric integral.
  \item The Jacobian involves the linearised operator~$L$, which decomposes
    into Fourier modes for radial $u_0$.
\end{itemize}

The direct certified programme for the sharp constant is:
\begin{enumerate}[leftmargin=2em]
  \item Prove a priori bounds:
    $\norm{f}_{L^2}=1$, $\norm{f}_{C^m}\leq B_m$,
    $0<\Lambda_{\min}\leq\Lambda\leq C_{\mathrm{ST}}^6$.
  \item Use Lemma~\ref{lem:positive-reduction} to restrict the extremiser
    search to boxes containing a nonnegative representative.
  \item Fix gauges (phase, rotation, translation); in the nonnegative
    slice, Lemma~\ref{lem:positive-gauge} fixes the translation parameter.
  \item Mesh $S^1$; represent $f$ by interval boxes plus derivative bounds.
  \item For each $(f,\Lambda)$ box, enclose the residual
    $E^*(|Ef|^4Ef)(\omega_i)-\Lambda f(\omega_i)$.
  \item Discard or validate via interval Newton.
  \item Sort validated $\Lambda_j$ intervals; certify the ordering.
\end{enumerate}

% =======================================================================
\section{Summary: The Logical Dependency Graph}
\label{sec:dependency}
% =======================================================================

We summarise the dependencies and status of the full programme.

\bigskip
\begin{center}
\fbox{\parbox{0.9\textwidth}{
\textbf{Proved:}
\begin{itemize}[leftmargin=1.5em, topsep=2pt, itemsep=1pt]
  \item Spectral support forcing ($\wh u\subseteq S^1$)
  \item 1D case complete
  \item Compactness of linearised operator $K_{u_0}$ on $L^2(S^1)$
  \item Fredholm index zero for $L=K-\Lambda I$
  \item Eigenvalue decomposition at the constant solution
  \item Nondegeneracy $\Leftrightarrow$ Bessel integrals~\eqref{eq:nondeg-cond}
  \item $\mathcal J_n\to 0$ (finitely many integrals to check)
  \item Regularity bootstrap for critical points
  \item Inverse Strichartz inequality for $S^1$ extension
  \item Profiles of critical sequences are critical points
  \item Existence of a maximiser (Theorem~\ref{thm:maximiser})
  \item Positivity reduction for extremisers
    (Lemma~\ref{lem:positive-reduction})
  \item Finiteness of nonnegative critical values in compact positive
    windows (Theorem~\ref{thm:positive-values-finite})
\end{itemize}

\textbf{Next concrete steps:}
\begin{itemize}[leftmargin=1.5em, topsep=2pt, itemsep=1pt]
  \item Compute $\mathcal J_n$ for $n=2,3,\ldots,n_*$ by interval
    arithmetic; verify~\eqref{eq:nondeg-cond}
  \item In the extremiser search, work with the finite positive-window
    list from Theorem~\ref{thm:positive-values-finite}
  \item Quantify the gap around $\Lambda_0$ by certifying
    $\nu_0\geq0.4\Lambda_0$ and running the far-region interval search
  \item Prove the constant is the unique maximiser modulo $G$
    (this alone gives $\mathcal R$)
\end{itemize}

\textbf{Open (harder):}
\begin{itemize}[leftmargin=1.5em, topsep=2pt, itemsep=1pt]
  \item Rule out multi-profile splitting for non-maximising
    critical sequences (Theorem~\ref{thm:compactness-mod-G})
  \item Enumeration / uniqueness of non-constant shapes
\end{itemize}

\textbf{Conditional conclusion:}
\begin{itemize}[leftmargin=1.5em, topsep=2pt, itemsep=1pt]
  \item Compactness + nondegeneracy $\Longrightarrow$ finitely many
    critical values $\Lambda_j$
  \item $\mathcal R=\max_j\Lambda_j^{1/6}$
\end{itemize}
}}
\end{center}

\bigskip
\noindent\rule{\linewidth}{0.4pt}
\begin{center}
  \small\itshape
  Notes compiled from working discussions.  Statements marked \proven{}
  carry complete proofs; statements marked \conditional{} depend on
  explicitly stated hypotheses; statements marked \open{} require new
  analytic input.
\end{center}

\end{document}
